\documentclass[a4paper,12pt]{article}

\usepackage[T1]{fontenc}
\usepackage[italian]{babel}
\usepackage{graphicx}
\usepackage{geometry}
\usepackage{tabularx} % Per l'ambiente tabularx (tabelle)
\usepackage{calc} % Sempre per le tabelle
\usepackage[table]{xcolor}
\usepackage[most]{tcolorbox}
\usepackage{fancyhdr}
\usepackage[hidelinks]{hyperref}
\usepackage{tikz}
\usepackage{pgf-pie}
\usepackage[utf8]{inputenc}
\usepackage{booktabs}
\usepackage{xcolor}  
\usepackage{float}
\usepackage{hyperref}
\usepackage{titlesec}
\usepackage{siunitx}
\usepackage{nicematrix}
\usepackage{colortbl}
\usepackage{eurosym}
\usepackage[dvipsnames]{xcolor}
\usepackage{titlesec} 

\newcommand{\groupmail}{alittlebyte.swe@gmail.com}
\newcommand{\grouplogo}{../../.github/site-src/assets/logo_alittlebyte.png}
\newcommand{\groupsymbol}{../../.github/site-src/assets/solo_logo.png}

\geometry{
    left=2.5cm,
    right=2.5cm,
    top=2.5cm,
    bottom=2.5cm,
    headheight=331pt,
    includeheadfoot
}

\newcommand{\slice}[4]{
  \pgfmathparse{0.5*#1+0.5*#2}
  \let\midangle\pgfmathresult

  \draw[thick,fill=black!10] (0,0) -- (#1:1) arc (#1:#2:1) -- cycle;

  \node[label=\midangle:#4] at (\midangle:1) {};

  \pgfmathparse{min((#2-#1-10)/110*(-0.3),0)}
  \let\temp\pgfmathresult
  \pgfmathparse{max(\temp,-0.5) + 0.8}
  \let\innerpos\pgfmathresult
  \node at (\midangle:\innerpos) {#3};
}

\titleclass{\subsubsubsection}{straight}[\subsection]

\newcounter{subsubsubsection}[subsubsection]
\renewcommand\thesubsubsubsection{\thesubsubsection.\arabic{subsubsubsection}}
\renewcommand\theparagraph{\thesubsubsubsection.\arabic{paragraph}} % optional; useful if paragraphs are to be numbered

\titleformat{\subsubsubsection}
  {\normalfont\normalsize\bfseries}{\thesubsubsubsection}{1em}{}
\titlespacing*{\subsubsubsection}
{0pt}{3.25ex plus 1ex minus .2ex}{1.5ex plus .2ex}

\makeatletter
\renewcommand\paragraph{\@startsection{paragraph}{5}{\z@}%
  {3.25ex \@plus1ex \@minus.2ex}%
  {-1em}%
  {\normalfont\normalsize\bfseries}}
\renewcommand\subparagraph{\@startsection{subparagraph}{6}{\parindent}%
  {3.25ex \@plus1ex \@minus .2ex}%
  {-1em}%
  {\normalfont\normalsize\bfseries}}
\def\toclevel@subsubsubsection{4}
\def\toclevel@paragraph{5}
%\def\toclevel@paragraph{6}
\def\toclevel@subparagraph{6}
\def\l@subsubsubsection{\@dottedtocline{4}{7em}{4em}}
\def\l@paragraph{\@dottedtocline{5}{10em}{5em}}
\def\l@subparagraph{\@dottedtocline{6}{14em}{6em}}
\makeatother

\setcounter{secnumdepth}{5}
\setcounter{tocdepth}{5}



\pagestyle{fancy}
\fancyhf{}

\renewcommand{\headrulewidth}{0pt}
\fancyhead{}

\renewcommand{\footrulewidth}{0.4pt}
\fancyfoot[L]{\includegraphics[height=0.6cm]{\groupsymbol}\hspace{0.45em}aLittleByte}
\fancyfoot[R]{Piano di Progetto}

\fancypagestyle{firstpage}{
    \fancyhf{}
    \renewcommand{\headrulewidth}{0pt}
    \renewcommand{\footrulewidth}{0.4pt}
    \fancyhead[C]{%
        \makebox[\textwidth][c]{%
            \begin{tabular}{c}
                \includegraphics[height=10cm]{\grouplogo}\\[1ex]
                \small\texttt{\groupmail}
            \end{tabular}%
        }
    }
    \fancyfoot[L]{\includegraphics[height=0.6cm]{\groupsymbol}\hspace{0.45em}aLittleByte}
    \fancyfoot[R]{Piano di Progetto}
}

\providecommand{\glslink}[2]{\href{http://127.0.0.1:8765/glossary-html\#gls-#1}{\underline{#2}\textsuperscript{\scriptsize G}}}

\begin{document}

\thispagestyle{firstpage}

\vspace*{0.5cm}

\begin{center}
    {\LARGE \textbf{Piano di Progetto}}
\end{center}

\vspace{1.5cm}

\begin{center}
    \renewcommand{\arraystretch}{1.3}
    \begin{tcolorbox}[
        width=0.92\textwidth,
        colback=gray!6,
        colframe=gray!45,
        boxrule=0.5pt,
        arc=2mm,
        boxsep=0pt,
        left=8pt, right=8pt, top=8pt, bottom=8pt
    ]
        \begin{tabularx}{\linewidth}{>{\bfseries}p{3.5cm} X}
            Gruppo      & aLittleByte (gruppo 19) \\
            Versione & 1.0.0 \\
            Redattori   & A. Maule, A.Oloni, E. Bellet \\
            Verificatori &  V. Y. K. Fernandes, E. Bellet, D. Belluz \\
            Destinatari & Prof. Tullio Vardanega, Prof. Riccardo Cardin \\
            Data        & 28/05/2026
        \end{tabularx}
    \end{tcolorbox}
\end{center}

\newpage
\fancyfoot[R]{Piano di Progetto\hspace{0.8em}}
\newgeometry{left=2.5cm,right=2.5cm,top=2.5cm,bottom=2.5cm,includefoot}

\begin{center}
    {\Large \textbf{Registro delle Modifiche}}
\end{center}

\vspace{0.35cm}
\renewcommand{\arraystretch}{1.35}
\rowcolors{2}{gray!8}{white}
\begin{table}[h]
\centering
    \setlength{\tabcolsep}{3pt} % Riduce leggermente lo spazio tra le colonne
    \resizebox{\textwidth}{!}{% % Forza il ridimensionamento orizzontale
    \footnotesize
    \begin{tabularx}{\textwidth}{|l|l|X|p{2cm}|p{2.5cm}|p{2.5cm}|}
        \hline
        \rowcolor{gray!20}
        \textbf{Versione} & \textbf{Data} & \textbf{Descrizione} & \textbf{Redattore} & \textbf{Verificatore} & \textbf{Approvatore} \\
        \hline 
        1.0.0 & 28/05/2026 & Approvazione versione 1.0.0 & - & - & Angela Maule \\
        \hline
        0.9.0 & 24/05/2026 & Aggiunta della tabella ore rimanenti per ruolo & Eleonora Bellet & Diego Belluz & - \\
        \hline
        0.8.1 & 24/05/2026 & Completato il quarto sprint & Eleonora Bellet & Diego Belluz & - \\
        \hline
        0.8.0 & 16/05/2026 & Aggiunta del consuntivo quarto sprint & Eleonora Bellet & Diego Belluz & - \\
        \hline
        0.7.0 & 15/05/2026 & Aggiunti i preventivi a finire del terzo sprint e iniziato il preventivo del quarto sprint & Eleonora Bellet & Diego Belluz & - \\
        \hline
        0.6.1 & 11/05/2026 & Aggiunti i preventivi a finire del primo e del secondo sprint & Eleonora Bellet & Diego Belluz & - \\
        \hline
        0.6.0 & 09/05/2026 & Aggiunto terzo sprint & Alex Oloni & V. Y. K. Fernandes & - \\
        \hline
        0.5.1 & 05/05/2026 & Aggiunti rischi all'analisi dei rischi & Alex Oloni & V. Y. K. Fernandes & - \\
        \hline
        0.5.0 & 02/05/2026 & Aggiunta di specifiche nel Primo Periodo e Stesura del Secondo Periodo & Alex Oloni & Eleonora Bellet & - \\
        \hline
        0.4.0 & 29/04/2026 & Inizio stesura della sezione §4 & Alex Oloni & Eleonora Bellet & - \\
        \hline
        0.3.2 & 26/04/2026 & Stesura della sezione §3 & Angela Maule & V. Y. K. Fernandes & - \\
        \hline
        0.3.1 & 23/04/2026 & Stesura delle sezioni §1 e §2 & Angela Maule & V. Y. K. Fernandes & - \\
        \hline
        0.3.0 & 21/04/2026 & Aggiunta delle sezioni: §1.2, §1.3, §1.4, §2, §3, Preventivo ed aggiunta specifiche sulla sezione Introduzione e Pianificazione & Angela Maule & V. Y. K. Fernandes & - \\
        \hline
        0.2.0 & 14/04/2026 & Stesura del primo periodo & Angela Maule & V. Y. K. Fernandes & - \\
        \hline
        0.1.0 & 13/04/2026 & Prima stesura del documento & Angela Maule & V. Y. K. Fernandes & - \\
        \hline
    \end{tabularx}}
    \caption{Registro delle Modifiche}
\end{table}
\newpage
\pagenumbering{arabic}
\setcounter{page}{1}
\fancyfoot[R]{Piano di Progetto\hspace{0.8em}\thepage}
\newpage

\tableofcontents
\listoftables
\listoffigures

\newpage

\section{Introduzione}
\label{sec:introduzione}

\subsection{Scopo del documento}
\label{sec:scopo_del_documento}
Il \glslink{piano-di-progetto}{Piano di Progetto} è un documento nella quale vengono fissati in maniera chiara e coincisa gli obiettivi, la progettazione, l'esecuzione ed il controllo delle attività del progetto, fungendo quindi da bussola per allineare il lavoro in svolgimento verso un'unico obiettivo comune.
Inoltre questo documento specifica anche la suddivisione del lavoro, la stima dei costi di realizzazione, le risorse necessarie e i vari rischi con la loro conseguente \glslink{mitigazione-dei-rischi}{mitigazione}, dati necessari per il compimento dei traguardi stabiliti ed evitare costi ed espansioni non controllate.


\subsection{Scopo del prodotto}
\label{sec:scopo_del_prodotto}
Il prodotto da realizzare consiste in un'estensione innovativa per \textit{''\glslink{nexum}{Nexum}"}, una piattaforma digitale sviluppata dall'azienda \glslink{committente-proponente}{proponente} \textit{Eggon}, dedicata alla comunicazione interna e alla gestione HR.
Nello specifico, il progetto mira a potenziare l'ecosistema \glslink{nexum}{Nexum} attraverso lo sviluppo e l'integrazione di due moduli basati sull'\glslink{intelligenza-artificiale-ai}{Intelligenza Artificiale}:
\begin{itemize}
    \item \textbf{\glslink{ai-assistant-generativo}{AI Assistant Generativo}:} Un modulo dedicato agli utenti amministrativi (\glslink{dashboard}{dashboard}) che permette la creazione autonoma di contenuti aziendali (titolo, testo, immagini di copertina) a partire da un \glslink{prompt}{prompt}, garantendo l'adeguamento automatico del tono e dello stile del messaggio.
    \item \textbf{\glslink{ai-co-pilot-per-consulenti-del-lavoro-cdl}{AI Co-Pilot per Consulenti del Lavoro (CdL)}:} Un sistema avanzato in grado di riconoscere automaticamente la tipologia di documenti caricati (es. cedolini, comunicazioni), estrarne i destinatari tramite tecniche di \glslink{ocr-riconoscimento-ottico-dei-caratteri}{OCR} ed \textbf{\glslink{entity-resolution}{Entity Resolution}}, e automatizzare lo smistamento e il \glslink{dispaccio}{dispaccio} massivo.
\end{itemize}

L'obiettivo finale è migliorare l'efficienza dei processi HR e semplificare l'interoperabilità tra le aziende e gli studi dei Consulenti del Lavoro.

\subsection{Glossario}
\label{sec:glossario}
Per questioni di ordine e di chiarezza dei concetti utilizzati, si è deciso di redarre un \textit{\glslink{glossario}{Glossario}}, nel quale viene contenuta
tutta la terminologia e la specifica dei termini che possono risultare confusionari, tecnici o semplicemente ambigui.

Il \glslink{glossario}{glossario} è disponibile a questo link: da mettere

\subsection{Riferimenti}
\label{sec:riferimenti}
\subsubsection{Riferimenti normativi}
\begin{itemize}
    \item \glslink{norme-di-progetto}{Norme di Progetto}
    \item \glslink{capitolato}{Capitolato} C5: \href{https://www.math.unipd.it/~tullio/IS-1/2025/Progetto/C5.pdf}{\textit{\textbf{'Nexum'}}}
    \item \href{https://www.math.unipd.it/~tullio/IS-1/2025/Dispense/PD1.pdf}{Regolamento Progetto Didattico}
\end{itemize}
\subsubsection{Riferimenti informativi}
\begin{itemize}
    \item \glslink{glossario}{Glossario}
    \item \href{https://www.math.unipd.it/~tullio/IS-1/2025/Dispense/T02.pdf}{T2: Ciclo di Vita del Software}
    \item \href{https://www.math.unipd.it/~tullio/IS-1/2025/Dispense/T03.pdf}{T3: Modelli di sviluppo Software}
    \item \href{https://www.math.unipd.it/~tullio/IS-1/2025/Dispense/T04.pdf}{T4: Gestione di Progetto}
\end{itemize}


\section{Analisi dei Rischi}
\label{sec:analisi_dei_rischi}
In questa sezione vengono individuati e analizzati i fattori di \glslink{rischio}{rischio} che, a seconda del loro livello di gravità, potrebbero ostacolare l'avanzamento dei lavori e il successo dell'iniziativa. Rilevare tempestivamente questi rischi e predisporre le opportune tecniche di \glslink{mitigazione-dei-rischi}{mitigazione} permette di ridurre drasticamente gli imprevisti in fase di sviluppo. Questa gestione controllata favorisce il rispetto del budget preventivato, l'aderenza alle tempistiche di consegna e una maggiore fluidità nei processi di lavoro. Il processo di gestione \glslink{mitigazione-dei-rischi}{dei rischi} è strutturato nelle seguenti fasi sequenziali:
\begin{itemize}
    \item \textbf{Identificazione}: Questa fase consiste nell'individuare le varie zone di \glslink{rischio}{rischio}, quale area di competenza viene impattata ed il fattore scatenante.
    In esso viene creato un elenco/inventario per far si che venga creata una mappatura completa sui vari rischi che possono compromettere il progetto.
    \item \textbf{Analisi}: Valutazione dettagliata e approfondita \glslink{mitigazione-dei-rischi}{dei rischi}, per valutarne la probabilità e la gravità che ne consegue se il \glslink{rischio}{rischio} si applica.
    Analizzando i vari rischi, si può così optare per strategie più efficaci e funzionali per il trattamento di essi.
    \item \textbf{Valutazione}: Consiste nella specifica delle criticità, in ordine di importanza, per poi proporre la \glslink{mitigazione-dei-rischi}{mitigazione} necessaria e più adatta per quel determinato \glslink{rischio}{rischio},
    Grazie a ciò, si ha un'idea più chiara su quali minacce concentrarsi maggiormente, ottimizzando i tempi e l'uso di risorse.
    \item \textbf{Gestione}: Implica l'utilizzo di mezzi e tecniche come misure preventive per evitare l'insorgere \glslink{mitigazione-dei-rischi}{dei rischi}, utilizzando proprio le mitigazioni.
    Quindi si specifica che, dopo le fasi precedenti, si procede ad eseguire azioni concrete con l'obiettivo di bloccare i vari contrattempi che possono nascere.
    \item \textbf{Monitoraggio e Revisione}: Dopo aver implementato le varie mitigazioni, il periodo che segue consiste nel controllo periodico e regolare delle soluzioni proposte e messe in atto,
    per far si che le soluzioni adottate siano efficienti e coprano le problematiche che nascono, oltre a modificare la strategia se nuove difficoltà si presentano.
\end{itemize}

\subsection{Categorie di Rischio}
Queste fasi elencate devono essere attuate costantemente durante tutto il periodo di vita del progetto, con lo scopo di avere efficenza a livello di tempi e di costi.
Per identificare meglio questi rischi, si è deciso di utilizzare un formato nel quale viene specificata la zona di interesse e l'importanza del \glslink{rischio}{rischio}:
\begin{center}
    \textbf{R[Tipo][Indice]}
\end{center}
Dove \textbf{Tipo} indica la categoria in cui rientra il \glslink{rischio}{rischio}, che si divide in tre sezioni:
\begin{itemize}
    \item \textbf{T}: Tecnologico
    \item \textbf{O}: Organizzativo, legato alla comunicazione e alla gestione interna del team
    \item \textbf{P}: Personale, singolare ad ogni membro del gruppo
    \item \textbf{R}: Requisiti, legato a possibili modifiche, aggiunte o incomprensioni delle richieste del \glslink{committente-proponente}{proponente}
\end{itemize}
L'\textbf{\textit{Indice}} indica l'ordine \glslink{mitigazione-dei-rischi}{dei rischi} su ogni categoria, in cui abbiamo deciso di metterli in ordine di Probabilità e Pericolosità.

\newpage
\subsection{RT: Rischi legati alle Tecnologie}
\subsubsection{RT1: Complessità nell'apprendimento delle nuove tecnologie}
Nel \glslink{capitolato}{Capitolato} C5 \textbf{\textit{'\glslink{nexum}{Nexum}'}} richiede varie tecnologie richieste dall'azienda \glslink{committente-proponente}{proponente} che escono dal nostro campo di competenza che abbiamo sviluppato durante il periodo di Università, andando ad imparare tecnologie di cui i membri non hanno \glslink{confidenza-confidence-score}{confidenza} e padronanza.

\vspace{0.5cm}
\hrule
\vspace{0.5em}
\textbf{RT1: Complessità nell'apprendimento delle nuove tecnologie}
\vspace{0.5em}
\hrule
\vspace{1cm}

\begin{table}[h]
\begin{tabular}{>{\bfseries}p{3cm} | p{12.2cm}}
Codice & RT1 \\

Probabilità & Molto Alta\\

Pericolosità & Alta\\

Conseguenze & L'apprendimento e l'utilizzo con minima sicurezza dei nuovi linguaggi e tecnologie proposte, comporterà sicuramente ad un rallentamento nel raggiungimento degli obiettivi del progetto,
portando tutto il gruppo a concentrarsi ed adattarsi alle nuove risorse proposte. \\

\glslink{mitigazione-dei-rischi}{Mitigazione} & È indispensabile che ogni membro, sia singolarmente che a gruppi, dedichi un periodo di tempo per imparare ed esplorare le nuove tecnologie, per far si che vengano utilizzate nel modo più efficiente possibile all'interno del \glslink{capitolato}{capitolato}.
Ogni membro del gruppo si dedicherà nella ricerca di varia documentazione per l'apprendimento delle nuove tecnologie e le varie best practices.\\
\end{tabular}
\caption{RT1: Complessità nell'apprendimento delle nuove tecnologie}
\end{table}
\newpage

\subsubsection{RT2: Mancanza o insufficienza di documenti o risorse per le tecnologie}
Avendo già difficoltà non indifferenti con l'apprendimento delle nuove tecnologie, se le risorse e la documentazione risultano insufficienti o carenti, il processo di studio diventa ancora più complesso, principalmente per le tecnologie non note a nessun membro del team, oltre ad adattarle al progetto.
\vspace{0.5cm}
\hrule
\vspace{0.5em}
\textbf{RT2: Mancanza o insufficienza di documenti o risorse per le tecnologie}
\vspace{0.5em}
\hrule
\vspace{1cm}
\begin{table}[h]
\begin{tabular}{>{\bfseries}p{3cm} | p{12.2cm}}
Codice & RT2 \\

Probabilità & Alta\\

Pericolosità & Alta\\

Conseguenze &  Senza documentazione adatta o carente, questo \glslink{rischio}{rischio} causa rallentamenti, per l'apprendimento e la comprensione dei linguaggi, e ritardi nello sviluppo del codice e nella fase di debugging \\

\glslink{mitigazione-dei-rischi}{Mitigazione} & Si può richiedere all'azienda \glslink{committente-proponente}{proponente} di proporci qualche spunto o documentazione necessaria per imparare le nuove tecnologie, oppure si possono
trovare dei compromessi sull'utilizzo delle tecnologie, utilizzando delle soluzioni alternative meno complicate e con la documentazione più completa possibile\\
\end{tabular}
\caption{RT2: Mancanza o insufficienza di documenti o risorse per le tecnologie}
\end{table}
\newpage

\subsubsection{RT2: Disallineamento e integrazione dei sottosistemi}
Nello sviluppo di un sistema complesso come quello richiesto per il \glslink{capitolato}{Capitolato} C5 \textbf{\textit{'\glslink{nexum}{Nexum}'}}, il lavoro viene diviso in moduli (frontend, backend, database). Esiste il \glslink{rischio}{rischio} concreto che le interfacce di comunicazione tra questi sottosistemi non siano perfettamente allineate, causando errori critici in fase di unione del codice.

\vspace{0.5cm}
\hrule
\vspace{0.5em}
\textbf{RT3: Disallineamento e integrazione dei sottosistemi}
\vspace{0.5em}
\hrule
\vspace{1cm}

\begin{table}[h]
\begin{tabular}{>{\bfseries}p{3cm} | p{12.2cm}}
Codice & RT3 \\

Probabilità & Media \\

Pericolosità & Alta \\

Conseguenze & Il sistema potrebbe risultare frammentato e inutilizzabile nonostante il corretto funzionamento dei singoli componenti presi isolatamente, costringendo il team a pesanti sessioni di refactoring non pianificate. \\

\glslink{mitigazione-dei-rischi}{Mitigazione} & Definizione formale e preventiva dei contratti API e dei protocolli di comunicazione tra i vari moduli. Introduzione di \glslink{test-di-integrazione}{test di integrazione} automatici all'interno del flusso di Continuous Integration. \\
\end{tabular}
\caption{RT3: Disallineamento e integrazione dei sottosistemi}
\end{table}
\newpage

\subsubsection{RT4: Incongruenze tra ambienti di sviluppo e produzione}
La varietà di sistemi operativi e configurazioni hardware tra i sette membri del gruppo può portare a bug difficili da replicare, noti come problemi di "ambiente". Questo \glslink{rischio}{rischio} può causare il malfunzionamento del software sulla macchina del \glslink{committente-proponente}{proponente} nonostante funzioni sui PC del team.

\vspace{0.5cm}
\hrule
\vspace{0.5em}
\textbf{RT4: Incongruenze tra ambienti di sviluppo e produzione}
\vspace{0.5em}
\hrule
\vspace{1cm}

\begin{table}[h]
\begin{tabular}{>{\bfseries}p{3cm} | p{12.2cm}}
Codice & RT4 \\

Probabilità & Media \\

Pericolosità & Media \\

Conseguenze & Rallentamento delle attività di debugging e \glslink{rischio}{rischio} di fallimento delle dimostrazioni ufficiali a causa di dipendenze mancanti o configurazioni errate nel sistema di destinazione. \\

\glslink{mitigazione-dei-rischi}{Mitigazione} & Standardizzazione dell'ambiente di sviluppo tramite l'utilizzo di containerizzazione (\textit{Docker}). Ogni membro del gruppo utilizzerà la medesima immagine per garantire che il software si comporti in modo identico su ogni macchina. \\
\end{tabular}

\caption{RT4: Incongruenze tra ambienti di sviluppo e produzione}
\end{table}
\newpage 

\subsubsection{RT5: Esposizione involontaria di dati sensibili o credenziali}
Durante il caricamento del codice su piattaforme di versionamento pubblico come \glslink{github}{GitHub}, esiste il \glslink{rischio}{rischio} di includere accidentalmente chiavi API, password di database o token di accesso. Questo metterebbe a \glslink{rischio}{rischio} la sicurezza del progetto.

\vspace{0.5cm}
\hrule
\vspace{0.5em}
\textbf{RT5: Esposizione involontaria di dati sensibili o credenziali}
\vspace{0.5em}
\hrule
\vspace{1cm}

\begin{table}[h]
\begin{tabular}{>{\bfseries}p{3cm} | p{12.2cm}}
Codice & RT5 \\

Probabilità & Bassa \\

Pericolosità & Molto Alta \\

Conseguenze & Compromissione degli account, possibile utilizzo non autorizzato delle risorse cloud da parte di terzi e necessità di invalidare e rigenerare tutte le credenziali esposte, con relativa pulizia dello storico \glslink{git}{Git}. \\

\glslink{mitigazione-dei-rischi}{Mitigazione} & Utilizzo sistematico di file di configurazione locali (.env) esclusi dal tracciamento \glslink{git}{Git} tramite il file .gitignore. Formazione dei membri sull'uso di segreti e adozione di tool automatici per la scansione dei segreti nei commit. \\
\end{tabular}
\caption{RT5: Esposizione involontaria di dati sensibili o credenziali}
\end{table}
\newpage

\subsubsection{RT6: Perdita di file o danneggiamento di essi}
C'è la possibilità che i file vengano danneggiati o vengono cancellati per errori software, hardware o per errori umani.
\vspace{0.5cm}
\hrule
\vspace{0.5em}
\textbf{RT6: Perdita di file o danneggiamento di essi}
\vspace{0.5em}
\hrule
\vspace{1cm}

\begin{table}[h]
\begin{tabular}{>{\bfseries}p{3cm} | p{12.2cm}}
Codice & RT6 \\

Probabilità & Bassa\\

Pericolosità & Media\\

Conseguenze & Perdere i documenti necessari causa un rallentamento non oscurabile al proseguimento del progetto, dovendo rimediare all'errore con la ricostruzione dei file mancanti\\

\glslink{mitigazione-dei-rischi}{Mitigazione} & Si adotta una soluzione nella quale viene utilizzato un sistema per il salvataggio dei file robusto ed affidabile, con salvataggio del lavoro graduale ed incrementale. 
Fare ciò garantisce sicurezza ed integrità dei documenti di lavoro, facilitandone anche il recupero in caso di perdita. \\
\end{tabular}
\caption{RT6: Perdita di file o danneggiamento di essi}
\end{table}
\newpage

\subsection{RO: Rischi legati all'Organizzazione del gruppo}
\subsubsection{RO1: Mancanza di comunicazione chiara e completa}
La comunicazione all'interno del gruppo può essere insufficientemente esaustiva o chiara, portando ad errori di comprensione, malintesi 
e disallineamenti sulle attività pianificate

\vspace{0.5cm}
\hrule
\vspace{0.5em}
\textbf{RO1: Mancanza di comunicazione chiara e completa}
\vspace{0.5em}
\hrule
\vspace{1cm}

\begin{table}[h]
\begin{tabular}{>{\bfseries}p{3cm} | p{12.2cm}}
Codice & RO1 \\

Probabilità & Alta\\

Pericolosità & Alta\\

Conseguenze & Può causare gravi errori nello sviluppo, portare a rallentamenti e ritardi consistenti per rimediare alle mancanze o duplicazioni delle attività.
Oltre al malumore generale causato nel gruppo, comporta a probabili errori nel prodotto finale, peggiorando la qualità dello stesso.\\

\glslink{mitigazione-dei-rischi}{Mitigazione} & Si pubblica un resoconto chiaro e coinciso delle mansioni che si devono portare a termine, ognuno col proprio destinatario, e discutendo in maniera costante con gli altri membri attraverso i vari canali di comunicazione (WhatsApp, Telegram, Discord). Si organizzano riunioni periodiche per portare agli altri membri il riassunto dettagliato del proprio lavoro durante il periodo di \glslink{sprint}{sprint}, discutendo inoltre delle prossime \glslink{milestone}{milestone} che si vogliono raggiungere, fissando poi in maniera esaustiva le varie mansioni di ogni singolo membro.\\
\end{tabular}
\caption{RO1: Mancanza di comunicazione chiara e completa}
\end{table}
\newpage


\subsubsection{RO2: Confusione sulla gestione delle responsabilità}
La mancata chiarezza sui ruoli affidati durante il periodo può portare a confusione sulle proprie attività, sovrapposizioni e lacune in settori che richiedono maggiore attenzione.
\vspace{0.5cm}
\hrule
\vspace{0.5em}
\textbf{RO2: Confusione sulla gestione delle responsabilità}
\vspace{0.5em}
\hrule
\vspace{1cm}

\begin{table}[h]
\begin{tabular}{>{\bfseries}p{3cm} | p{12.2cm}}
Codice & RO2 \\

Probabilità & Alta\\

Pericolosità & Alta\\

Conseguenze & La mancata organizzazione sulla divisione delle responsabilità può portare a ritardi e ad uno spreco di risorse non indifferente, con il \glslink{rischio}{rischio} di 
far slittare ulteriormente la scadenza della consegna del progetto\\

\glslink{mitigazione-dei-rischi}{Mitigazione} & Durante le riunioni del team, si definiscono in maniera precisa e chiara, scrivendo le varie responsabilità e attività da svolgere durante i singoli \glslink{sprint}{Sprint}, così da evitare fraintendimenti e duplice lavoro.\\
\end{tabular}
\caption{RO2: Confusione sulla gestione delle responsabilità}
\end{table}
\newpage

\subsubsection{RO3: Gestione del tempo e delle scadenze}
La mancanza di organizzazione del tempo e delle varie scadenze indicate per il compimento del progetto, come la sottovalutazione degli impegni o scadenze irrealizzabili, possono portare a ritardi che
possono compromettere il risultato finale, oltre a causare slittamenti nelle scadenze.
\vspace{0.5cm}
\hrule
\vspace{0.5em}
\textbf{RO3: Gestione del tempo e delle scadenze}
\vspace{0.5em}
\hrule
\vspace{1cm}

\begin{table}[h]
\begin{tabular}{>{\bfseries}p{3cm} | p{12.2cm}}
Codice & RO3 \\

Probabilità & Media\\

Pericolosità & Alta\\

Conseguenze & La sottovalutazione dei propri impegni può portare a rimandare il lavoro a cui si è stati affidati, causando slittamenti non indifferenti riguardo le scadenze.
Inoltre la proposizione di pianificazioni irrealizzabili, portano a svolgere le attività in maniera sbrigativa, portando ad una qualità peggiore del prodotto finale. \\

\glslink{mitigazione-dei-rischi}{Mitigazione} & Ogni membro del gruppo deve organizzarsi affinchè i vari compiti che gli vengono assegnati siano compatibili e siano realizzabili. Il \glslink{responsabile}{Responsabile} deve
essere inoltre avvisato tempestivamente se ci sono difficoltà nel rispettare le scadenze, cercando di darsi una mano in maniera reciproca. Infine, devono essere svolte riunioni periodiche, affinchè ogni difficoltà emersa possa essere resa nota a tutti i membri e così adattarsi di conseguenza.\\
\end{tabular}
\caption{RO3: Gestione del tempo e delle scadenze}
\end{table}
\newpage

\subsubsection{RO4: Esaurimento precoce del monte ore per un ruolo specifico}
A causa dei rigidi vincoli di ripartizione oraria stabiliti nel \glslink{preventivo}{preventivo} e della regola di rotazione, esiste la possibilità che un membro del team esaurisca prematuramente le ore a sua disposizione per un determinato ruolo (ad esempio \glslink{amministratore}{Amministratore} o \glslink{progettista}{Progettista}), prima del termine del progetto.

\vspace{0.5cm}
\hrule
\vspace{0.5em}
\textbf{RO4: Esaurimento precoce del monte ore per un ruolo specifico}
\vspace{0.5em}
\hrule
\vspace{1cm}

\begin{table}[h]
\begin{tabular}{>{\bfseries}p{3cm} | p{12.2cm}}
Codice & RO4 \\

Probabilità & Media\\

Pericolosità & Media\\

Conseguenze & Impossibilità formale per un componente di svolgere determinate mansioni senza violare le direttive del \glslink{piano-di-progetto}{Piano di Progetto}. Necessità di riassegnare i task in emergenza ad altri membri, con possibile rallentamento delle attività in corso e perdita di efficienza.\\

\glslink{mitigazione-dei-rischi}{Mitigazione} & Monitoraggio costante, alla fine di ogni \glslink{sprint}{Sprint}, del bilancio orario rimanente (Saldo Ore) per ciascun componente tramite il foglio di calcolo condiviso. Quando un membro si avvicina all'esaurimento delle ore per un ruolo, il \glslink{responsabile}{Responsabile} pianificherà in anticipo un passaggio di consegne verso un collega con un monte ore ancora intatto.\\
\end{tabular}
\caption{RO4: Esaurimento precoce del monte ore per un ruolo specifico}
\end{table}
\newpage

\subsubsection{RO5: Mancata o carenza di comunicazione col proponente}
La comunicazione con l'azienda \glslink{committente-proponente}{proponente} è fondamentale per far si che il prodotto finale sia in linea con le richieste, se questa viene a mancare o le risposte non sono abbastanza tempestive, possono causare ritardi e fraintendimenti su questioni tecnologiche ed organizzative.
\vspace{0.5cm}
\hrule
\vspace{0.5em}
\textbf{RO5: Mancata o carenza di comunicazione col \glslink{committente-proponente}{proponente}}
\vspace{0.5em}
\hrule
\vspace{1cm}

\begin{table}[h]
\begin{tabular}{>{\bfseries}p{3cm} | p{12.2cm}}
Codice & RO5 \\

Probabilità & Media\\

Pericolosità & Media\\

Conseguenze & Ritardi a causa di incertezze su domande a livello tecnico ed organizzativo. Probabile disallineamento tra la richiesta del \glslink{committente-proponente}{proponente}
ed il prodotto del team\\
\glslink{mitigazione-dei-rischi}{Mitigazione} & Ogni qualvolta ci sono dubbi di rilievo, è opportuno interpellare il \glslink{committente-proponente}{proponente} tramite i mezzi di comunicazione da loro fornitoci. Se il problema persiste,
si contatta il \glslink{committente-proponente}{committente} per un colloquio. È inoltre importante fare una documentazione chiara e precisa sui verbali esterni.\\
\end{tabular}
\caption{RO5: Mancata o carenza di comunicazione col proponente}
\end{table}

\newpage

\subsection{RP: Rischi personali legati ai singoli membri}
\subsubsection{RP1: Mancato o insufficiente svolgimento dei propri impegni}
Il mancato svolgimento dei propri impegni per problemi di salute, accademici o personali possono compromettere la continuità del progetto.
\vspace{0.5cm}
\hrule
\vspace{0.5em}
\textbf{RP1: Mancato o insufficiente svolgimento dei propri impegni}
\vspace{0.5em}
\hrule
\vspace{1cm}

\begin{table}[h]
\begin{tabular}{>{\bfseries}p{3cm} | p{12.2cm}}
Codice & RP1 \\

Probabilità & Alta\\

Pericolosità & Alta\\

Conseguenze & La riduzione del tempo impiegato al progetto porta ad un rallentamento delle varie attività. Procrastinazione e mancanza di organizzazione dei propri impegni,
i quali comportano problemi nella \glslink{pianificazione}{pianificazione} delle varie attività. Eventi imprevisti come malattia o emergenze, oppure situazioni lavorative per alcuni membri, sicuramente 
incidono notevolmente sul tempo di termine del progetto, allungando i tempi delle scadenze ed influenzando le capacità del gruppo.\\

\glslink{mitigazione-dei-rischi}{Mitigazione} & Pianificare in maniera chiara e realizzabile i propri impegni sul progetto, mettendoli in linea con i propri impegni personali ed accademici.
Prevedere del tempo extra in caso di eventi imprevisti. Comunicare in maniera tempestiva o, in caso di impegni accademici e lavorativi, a priori con un buon margine di tempo, al \glslink{responsabile}{Responsabile} del Progetto, affinchè si abbia maggior margine
di manovra per una miglior disposizione delle responsabilità e di redistribuzione delle attività \glslink{intelligenza-artificiale-ai}{ai} membri del gruppo, con effetto immediato o nel prossimo \glslink{sprint}{sprint}, in caso in cui 
non si vuole sovraccaricare di lavoro il team. \\
\end{tabular}
\caption{RP1: Mancato o insufficiente svolgimento dei propri impegni}
\end{table}

\newpage

\subsubsection{RP2: Lavoro svolto non conforme rispetto agli impegni dichiarati}
Se viene svolto un lavoro che non coincide con le specifiche del \glslink{capitolato}{capitolato} oppure non si adempie al lavoro svolto, ci sono \glslink{mitigazione-dei-rischi}{dei rischi} sul tempo e sui costi.
\vspace{0.5cm}
\hrule
\vspace{0.5em}
\textbf{RP2: Lavoro svolto non conforme rispetto agli impegni dichiarati}
\vspace{0.5em}
\hrule
\vspace{1cm}

\begin{table}[h]

\begin{tabular}{>{\bfseries}p{3cm} | p{12.2cm}}
Codice & RP2 \\

Probabilità & Media\\

Pericolosità & Alta\\

Conseguenze & Perdita di interesse del \glslink{committente-proponente}{proponente}, costi e tempi maggiori rispetto a quelli dichiarati per rimediare al problema creatosi e malumori
all'interno del gruppo, perdendo così coesione fra i vari membri del gruppo.\\

\glslink{mitigazione-dei-rischi}{Mitigazione} & Comunicazione continua fra i membri del gruppo. Adozione di una gestione chiara e rigorosa sul progetto, per evitare sbavature. Comunicazione periodica col cliente, affinchè il lavoro
sia in linea con gli impegni dichiarati\\
\end{tabular}
\caption{RP2: Lavoro svolto non conforme rispetto agli impegni dichiarati}
\end{table}

\newpage
\subsection{RR: Rischi legati ai Requisiti}
\subsubsection{RR1: Cambiamento o aggiunta di requisiti in corso d'opera}
Durante lo sviluppo del progetto, l'azienda \glslink{committente-proponente}{proponente} potrebbe modificare le proprie necessità o richiedere l'aggiunta di nuove funzionalità non previste inizialmente nel \glslink{capitolato}{capitolato} o nell'\glslink{analisi-dei-requisiti-adr}{Analisi dei Requisiti}.

\vspace{0.5cm}
\hrule
\vspace{0.5em}
\textbf{RR1: Cambiamento o aggiunta di requisiti in corso d'opera}
\vspace{0.5em}
\hrule
\vspace{1cm}

\begin{table}[h]
\begin{tabular}{>{\bfseries}p{3cm} | p{12.2cm}}
Codice & RR1 \\

Probabilità & Media\\

Pericolosità & Alta\\

Conseguenze & L'inserimento di nuovi requisiti ad architettura già definita comporta un pesante lavoro di refactoring, un aumento dei costi preventivati e un inevitabile ritardo nelle scadenze (\glslink{rtb-requirements-and-technology-baseline}{RTB} o \glslink{pb-product-baseline}{PB}), compromettendo la stabilità del lavoro già svolto.\\

\glslink{mitigazione-dei-rischi}{Mitigazione} & Adozione rigorosa della metodologia \glslink{agile}{Agile}. I nuovi requisiti verranno valutati dal team e, se accettati, inseriti nel Product Backlog per essere sviluppati negli \glslink{sprint}{Sprint} successivi, rinegoziando eventualmente i tempi di consegna con il \glslink{committente-proponente}{proponente}. Verrà data priorità \glslink{intelligenza-artificiale-ai}{ai} requisiti obbligatori per garantire in ogni caso la consegna dell'\glslink{mvp-minimum-viable-product}{MVP (Minimum Viable Product)}.\\
\end{tabular}
\caption{RR1: Cambiamento o aggiunta di requisiti in corso d'opera}
\end{table}
\newpage

\subsubsection{RR2: Ambiguità o errata interpretazione dei requisiti}
I requisiti concordati potrebbero risultare poco chiari, ambigui o interpretabili in modi diversi dai vari membri del team o tra il team e la \glslink{committente-proponente}{proponente} stessa.

\vspace{0.5cm}
\hrule
\vspace{0.5em}
\textbf{RR2: Ambiguità o errata interpretazione dei requisiti}
\vspace{0.5em}
\hrule
\vspace{1cm}

\begin{table}[h]
\begin{tabular}{>{\bfseries}p{3cm} | p{12.2cm}}
Codice & RR2 \\

Probabilità & Bassa\\

Pericolosità & Alta\\

Conseguenze & Sviluppo di funzionalità non conformi alle reali aspettative dell'azienda. Questo errore porta a un enorme spreco di risorse, in quanto il codice prodotto risulta inutile e deve essere scartato o pesantemente riscritto, generando frustrazione nel team.\\

\glslink{mitigazione-dei-rischi}{Mitigazione} & Durante la stesura del documento di \glslink{analisi-dei-requisiti-adr}{Analisi dei Requisiti}, ogni \glslink{caso-d-uso-uc}{caso d'uso} deve essere tracciato e approvato meticolosamente. Mantenere un dialogo costante con la \glslink{committente-proponente}{proponente}, inviando diagrammi e prototipi grafici (\glslink{mockup}{Mockup}) per validare visivamente le funzionalità prima di iniziare a scrivere il codice.\\
\end{tabular}
\caption{RR2: Ambiguità o errata interpretazione dei requisiti}
\end{table}
\newpage

\section{Stima Costi di Sviluppo}
In questa sezione viene dettagliato il budget preventivato per la realizzazione dell'intero progetto. Tale \glslink{pianificazione}{pianificazione} è stata redatta per consentire una gestione flessibile delle risorse umane, permettendo di intervenire sull'assegnazione delle ore per i vari ruoli qualora le necessità di sviluppo lo richiedano. Resta fermo l'impegno del gruppo a non superare il budget complessivo definito nel \glslink{preventivo}{Preventivo} dei Costi consegnato durante la fase di Candidatura.
\subsection{Preventivo Costi - Prima Stesura 24-03-2026}
\begin{table}[h!]
\centering
\renewcommand{\arraystretch}{1.5}
\begin{tabular}{|l|c|c|c|}
\hline
\textbf{Ruolo} & \textbf{Costo Orario (€/h)} & \textbf{Ore} & \textbf{Costo (€)} \\
\hline
\glslink{responsabile}{Responsabile} & 30€/h &56h  &1.680€  \\
\hline
\glslink{amministratore}{Amministratore} & 20€/h &51h  &1.020€  \\
\hline
\glslink{analista}{Analista} & 25€/h &60h  &1.500€  \\
\hline
\glslink{progettista}{Progettista} & 25€/h &143h  &3.575€  \\
\hline
\glslink{programmatore}{Programmatore} & 15€/h &160h  &2.400€  \\
\hline
\glslink{verificatore}{Verificatore} & 15€/h &160h  &2.400€  \\
\hline
\textbf{Totale} & & \textbf{630} & \textbf{12.575€} \\
\hline
\end{tabular}
\caption{Riassunto dei costi derivanti dalle ore assegnate a ciascun ruolo}
\label{tab:costi_ruoli}
\end{table}
\begin{figure}[h]
    \centering
    \begin{tikzpicture}
        \pie[
            text=legend,
            radius=3,
            color={red!60, blue!60, yellow!60, green!60, purple!60, orange!60}
        ]{
            8.9/\glslink{responsabile}{Responsabile} (56 ore),
            8.1/\glslink{amministratore}{Amministratore} (51 ore),
            9.5/\glslink{analista}{Analista} (60 ore),
            22.7/\glslink{progettista}{Progettista} (143 ore),
            25.4/\glslink{programmatore}{Programmatore} (160 ore),
            25.4/\glslink{verificatore}{Verificatore} (160 ore)
        }
    \end{tikzpicture}
    \caption{Preventivo Costi - Prima Stesura}
\end{figure}

\subsection{Ripartizione Oraria per Membro}
Come già previsto nella dichiarazione degli impegni, al fine di garantire il rispetto del principio di rotazione dei ruoli e assicurare un carico di lavoro uniforme, il team ha definito una matrice di ripartizione oraria. Ciascun componente del gruppo si impegna formalmente a prestare un totale di 90 ore produttive individuali lungo l'intero ciclo di vita del progetto, per un ammontare complessivo di 630 ore lavorative.

\begin{table}[H]
    \centering
    \renewcommand{\arraystretch}{1.5} 
    \begin{tabularx}{\textwidth}{|X|c|c|c|c|c|c|c|}
    \hline
    \textbf{Membro del Team} & \textbf{Resp.} & \textbf{Amm.} & \textbf{Ana.} & \textbf{Proget.} & \textbf{Progr.} & \textbf{Ver.} & \textbf{Totale} \\
    \hline
    Angelica Gastaldello & 8 & 7 & 9 & 21 & 22 & 23 & \textbf{90} \\
    \hline
    Eleonora Bellet & 8 & 7 & 8 & 20 & 23 & 24 & \textbf{90} \\
    \hline
    Alex Oloni & 8 & 8 & 8 & 20 & 23 & 23 & \textbf{90} \\
    \hline
    Diego Belluz & 8 & 7 & 9 & 21 & 23 & 22 & \textbf{90} \\
    \hline
    Angela Maule & 8 & 8 & 8 & 21 & 23 & 22 & \textbf{90} \\
    \hline
    Leonardo Soligo & 8 & 7 & 9 & 20 & 23 & 23 & \textbf{90} \\
    \hline
    V.Y. Kaneda Fernandes & 8 & 8 & 8 & 20 & 23 & 23 & \textbf{90} \\
    \hline
    \rowcolor{gray!20}
    \textbf{Totale Ore Ruolo} & \textbf{56} & \textbf{52} & \textbf{59} & \textbf{143} & \textbf{160} & \textbf{160} & \textbf{630} \\
    \hline
    \end{tabularx}
    \caption{Matrice di ripartizione oraria complessiva per ciascun componente del team}
    \label{tab:ripartizione_oraria_totale}
\end{table}
\newpage

\newpage

\section{Modello di Sviluppo}
\label{sec:modello_di_sviluppo}
Affinchè lo svolgimento del \glslink{capitolato}{Capitolato} \textit{C5: '\glslink{nexum}{Nexum}'} venga svolto in maniera efficiente ed in linea con le scadenze, il team ha deciso
di utilizzare un \glslink{modello-di-sviluppo}{modello di sviluppo} \href{https://agilemanifesto.org/}{\textbf{\textit{Agile}}} per la capacità di adattamento a cambiamenti improvvisi o naturali, oltre all'\glslink{approccio-incrementale}{approccio incrementale} ed iterativo, grazie al quale, con questi brevi cicli chiamati \textit{\glslink{sprint}{Sprint}},
offre una capacità maggiore di controllo e di miglioramento del software. Oltre a ciò, offre una continua contrattazione e dialogo col cliente, individuando problemi o modifiche necessarie in tempi brevi,
andando ad influire positivamente sul rispetto delle scadenze e dei costi di sviluppo.
\subsection{Fasi del Progetto}
Le fasi del progetto si dividono in tre macro sezioni, nelle quali ognuna di esse ha obiettivi specifici e diversificati:
\begin{itemize}
    \item \textbf{Candidatura}: Questo è il periodo preliminare e precedente all'inizio dell'\glslink{rtb-requirements-and-technology-baseline}{RTB}, nella quale vengono studiati e valutati i vari
    capitolati proposti dalle varie aziende proponenti.
    \item \textbf{\glslink{rtb-requirements-and-technology-baseline}{RTB} \textit{(\glslink{rtb-requirements-and-technology-baseline}{Requirements and Technology Baseline})}}: Questa fase avviene dopo l'aggiudicazione dell'appalto. Rappresenta il punto di controllo iniziale, fissato al termine della prima \glslink{milestone}{milestone} di revisione. Essa certifica il consolidamento dei requisiti, la stabilità della documentazione normativa e la validazione delle scelte tecnologiche avvenuta tramite il \textit{Proof of Concept}.
    \item \textbf{\glslink{pb-product-baseline}{PB}: \textit{(\glslink{pb-product-baseline}{Product Baseline})}}: In questo ultimo periodo, viene elaborato e sviluppato il prodotto software. È la baseline finale stabilita a conclusione della \glslink{milestone}{milestone} di accettazione, volta a validare la conformità dell'\glslink{mvp-minimum-viable-product}{MVP}, l'avvenuta esecuzione di tutti i piani di \glslink{test-di-sistema}{test di sistema} e la sottomissione finale del prodotto.
\end{itemize}

\subsection{Adozione del framework SCRUM al progetto}
Il team ha utilizzato SCRUM in questo progetto accademico per avere maggiore efficienza operativa ed un continuo riscontro con gli altri membri del gruppo.
Il ciclo di lavoro viene diviso in \glslink{sprint}{Sprint}, e in esso vengono eseguiti:
\begin{itemize}
    \item \textbf{\glslink{sprint}{Sprint} Planning}: All'inizio del nuovo \glslink{sprint}{Sprint}, il gruppo decide le varie mansioni e attività da svolgere emerse, producendo uno \textit{\glslink{sprint}{Sprint} Backlog}.
    \item \textbf{Meeting Sicrono}: Il team si riunisce in riunione su un'apposita applicazione in chiamata per discutere delle attività da svolgere
    durante lo \glslink{sprint}{Sprint} in apertura, specificando le varie priorità ed importanza, oltre alle risorse che andranno impiegate.
    \item \textbf{\glslink{sprint}{Sprint} Review e Retrospective}: Al termine dello \glslink{sprint}{Sprint}, si svolge una riunione sulla revisione dell'operato del team rispetto agli obiettivi
    fissati per il periodo, e si fanno delle considerazioni su vari aspetti per un miglioramento futuro della produttività.
\end{itemize}
\section{Pianificazione}
\label{sec:pianificazione}
\subsection{Calendario Revisioni del Progetto}
In questa sezione vengono illustrati le date previste per le revisioni del progetto e del suo conseguente avanzamento, definite in base alla quantità di lavoro da svolgere
e alle scadenze del corso.
\subsubsection{Date Avanzamento Progetto - Prima Stesura 29/04/2026}
\begin{table}[H]
    \centering
    \renewcommand{\arraystretch}{1.5} 
    \begin{tabularx}{\textwidth}{|X|X|}
    \hline \rowcolor{Dandelion!50}
    \textbf{Revisione} & \textbf{Data Stimata}\\
    \hline
    \textit{\glslink{rtb-requirements-and-technology-baseline}{Requirements and Technology Baseline} - \glslink{rtb-requirements-and-technology-baseline}{RTB}} & 2026/05/29\\
    \hline
    \textit{\glslink{pb-product-baseline}{Product Baseline} - \glslink{pb-product-baseline}{PB}} & 2026/09/10\\
    \hline
    \end{tabularx}
    \caption{Calendario Avanzamento Progetto - Prima Stesura}
\end{table}
\subsection{Periodo preliminare studio Capitolati (Pre-RTB)}
\begin{table}[H]
    \centering
    \renewcommand{\arraystretch}{1.5} % Aumenta l'altezza delle righe
    \begin{tabularx}{\textwidth}{|X|X|X|}
    \hline \rowcolor{Dandelion!50}
    \textbf{Data Inizio} & \textbf{Data Fine} & \textbf{Data Aggiudicazione Appalto}\\
    \hline
    06/03/2026 & 24/03/2026 & 26/03/2026\\
    \hline
    \end{tabularx}
    \caption{Periodo pre-RTB}
\end{table}

\subsubsection{Attività}
\label{sec:prertb_attivita}
Nella fase preliminare, con tutti i membri del team, si ha una discussione sull'opzione più interessante e fattibile a livello di tempi e costi.
Durante questa fase, il team prepara i documenti necessari per la candidatura e, conseguentemente, l'aggiudicazione dell'appalto, oltre ad un primo brainstorming sulle varie
task del \glslink{capitolato}{capitolato} scelto.
\subsubsection{Costruzione Candidatura}
\label{sec:costruzione_candidatura}
Nel periodo iniziale, il team \textit{aLittleByte} si è concentrato nella costruzione di una linea guida, per avere una coesione e comprensione reciproca sul lavoro da svolgere,
con l'obiettivo comune di eseguire un lavoro efficiente e preciso.
Nella fase a seguire, ogni membro ha contribuito nella stesura delle \textit{\glslink{norme-di-progetto}{Norme di Progetto}}, analizzando i vari capitolati proposti dalle aziende proponenti e condividendo
con tutti i propri interessi, dubbi ed incertezze su ogni \glslink{capitolato}{capitolato}. Dopo aver trovato i capitolati di maggior interesse e fattibilità, il team ha fissato dei colloqui con le aziende, per esporre le domande emerse durante le riunioni interne e maggiori specifiche organizzative e tecnologiche.
Parallelamente vengono svolte altre mansioni minori ma comunque essenziali, come la scelta del nome del gruppo, il logo, la creazione sia di una mail di riferimento che di una \glslink{repository}{repository} su
\glslink{github}{GitHub} dove andranno inseriti i vari file da presentare per la candidatura.

\subsection{Requirements and Technology Baseline - RTB}
\label{sec:rtb-requirements-and-technology-baseline}
\subsubsection{Obiettivi}
Durante questa fase, il team si impegna a redarre tutti i documenti necessari, ossia:
\begin{itemize}
    \item \textit{\glslink{norme-di-progetto}{Norme di Progetto}}
    \item \textit{\glslink{analisi-dei-requisiti-adr}{Analisi dei Requisiti}}
    \item \textit{\glslink{piano-di-progetto}{Piano di Progetto}}
    \item \textit{\glslink{piano-di-qualifica}{Piano di Qualifica}}
    \item \textit{\glslink{glossario}{Glossario}}
\end{itemize}
affinchè si abbia una chiara comprensione del \glslink{capitolato}{capitolato} scelto, e in certe situazioni si richiede all'azienda \glslink{committente-proponente}{proponente} incontri per un maggior riscontro sulle scelte attuate e su
eventuali compromessi tecnologici. Inoltre, tramite le scelte effettuate, verrà sviluppato un \textit{Proof of Concept}, ossia un prototipo funzionale in cui vengono inseriti i vari obiettivi indicati e le 
varie attività che deve svolgere, per valutarne quindi la complessità, fattibilità tecnologica e per avere un esempio tangibile delle varie soluzioni proposte.
Quindi i punti chiave in questa \textit{\glslink{rtb-requirements-and-technology-baseline}{Requirements and Technology Baseline}} sono essenzialmente:
\begin{itemize}
    \item \textbf{Stesura Documentazione}: Il gruppo si suddivide la creazione e la stesura della documentazione necessaria allo sviluppo del prodotto
    \item \textbf{Analisi}: Al \glslink{committente-proponente}{proponente} vengono richiesti incontri per un maggiore riscontro sullo sviluppo del Proof of Concept, affinchè si evitino ambiguità e si trovino compromessi o soluzioni migliori
    \item \textbf{Sviluppo del \textit{Proof of Concept}}: Viene quindi iniziato a sviluppare un prototipo del prodotto, per testarne funzionalità, attività e complessità.
\end{itemize}
\subsubsection{Periodi}
La fase \glslink{rtb-requirements-and-technology-baseline}{RTB} è suddivisa in 4 iterazioni (\glslink{sprint}{Sprint}).
\label{sec:rtb_primoperiodo}
\begin{table}[H]
    \centering
    \renewcommand{\arraystretch}{1.5} % Aumenta l'altezza delle righe
    \begin{tabularx}{\textwidth}{|X|X|X|}
    \hline \rowcolor{Dandelion!50}
    \textbf{\glslink{sprint}{Sprint}} & \textbf{Data Inizio} & \textbf{Data Fine}\\
    \hline
    \glslink{sprint}{Sprint} 1 & 30/03/2026 & 12/04/2026\\
    \hline
    \glslink{sprint}{Sprint} 2 & 13/04/2026 & 26/04/2026\\
    \hline
    \glslink{sprint}{Sprint} 3 & 27/04/2026 & 10/05/2026\\
    \hline
    \glslink{sprint}{Sprint} 4 & 11/05/2026 & 24/05/2026\\
    \hline

    \end{tabularx}
\end{table}

\subsubsection{Chiave di lettura dei periodi di Sprint}
Per garantire la massima trasparenza e agevolare la consultazione del documento, la rendicontazione di ciascuno \glslink{sprint}{Sprint} è stata strutturata secondo il seguente schema standard:

\begin{itemize}
    \item \textbf{Sintesi del periodo:} Un breve paragrafo introduttivo che riassume le principali attività affrontate dal team e gli obiettivi generali perseguiti durante l'iterazione.
    
    \item \textbf{Attività Pianificate (\glslink{sprint}{Sprint} Backlog):} Una tabella di dettaglio che elenca i task assegnati \glslink{intelligenza-artificiale-ai}{ai} vari ruoli, confrontando le ore preventivate con quelle effettivamente consumate a \glslink{consuntivo}{consuntivo}. Ogni attività è accompagnata dal relativo stato di avanzamento, indicato con le diciture ``\textit{Completata}'' o ``\textit{In corso}''.
    \begin{quote}
        \textit{Nota bene:} È fondamentale precisare che tali stati si riferiscono \textbf{esclusivamente agli incrementi e agli aggiornamenti pianificati per lo specifico \glslink{sprint}{Sprint}}, e non al completamento assoluto dell'intera attività o del documento a cui fanno riferimento. (A titolo di esempio, lo stato ``Completata'' per la voce ``Stesura \glslink{analisi-dei-requisiti-adr}{Analisi dei Requisiti}'' indica unicamente che il team ha portato a termine il carico di lavoro prefissato per quella specifica iterazione, non che il documento sia da considerarsi concluso o in versione definitiva).
    \end{quote}
    
    \item \textbf{Considerazioni (\glslink{sprint}{Sprint} Review e Retrospective):} Una sezione conclusiva in cui vengono analizzati criticamente i risultati raggiunti, evidenziando eventuali problematiche riscontrate durante lo sviluppo e formalizzando le azioni correttive adottate per le iterazioni successive.
\end{itemize}

\subsubsubsection{Primo Periodo}

\label{sec:rtb_primoperiodo}
\begin{table}[H]
    \centering
    \renewcommand{\arraystretch}{1.5} % Aumenta l'altezza delle righe
    \begin{tabularx}{\textwidth}{|X|X|}
    \hline \rowcolor{Dandelion!50}
    \textbf{Data Inizio} & \textbf{Data Fine}\\
    \hline
    30/03/2026 & 12/04/2026\\
    \hline

    \end{tabularx}
\end{table}


A seguito dell'aggiudicazione del \glslink{capitolato}{capitolato}, le attività del primo \glslink{sprint}{Sprint} si sono concentrate sull'analisi approfondita della documentazione e sul confronto con l'azienda \glslink{committente-proponente}{proponente} in merito alle scelte tecnologiche. Parallelamente, il team ha proceduto con l'individuazione dei \textit{casi d'uso} e ha avviato la stesura dell'\textit{\glslink{analisi-dei-requisiti-adr}{Analisi dei Requisiti}}, seguita dalle \textit{\glslink{norme-di-progetto}{Norme di Progetto}} e da una prima \glslink{bozza-draft}{bozza} del \textit{\glslink{glossario}{Glossario}}. Infine, è stata iniziata la realizzazione di un mock-up grafico, che fungerà da riferimento strutturale e visivo per le successive fasi di sviluppo.

\begin{center}
\noindent\textbf{Attivitá Pianificate \textit{(\glslink{sprint}{Sprint} Backlog)}}\\
\end{center}
Di seguito vengono dettagliate le singole attività pianificate per l'iterazione, con il relativo confronto tra le ore stimate in sede di \glslink{sprint}{Sprint} Planning e le ore effettivamente impiegate a \glslink{consuntivo}{consuntivo}.

\begin{table}[H]
    \centering
    \renewcommand{\arraystretch}{1.5}
    \resizebox{\textwidth}{!}{%
    \begin{tabular}{|l|c|c|c|c|}
    \hline
    \rowcolor{Dandelion!50}
    \textbf{Descrizione Incarico} & \textbf{Ruolo Assegnato} & \textbf{Ore Previste} & \textbf{Ore Effettive} & \textbf{Stato} \\
    \hline
    Aggiornamento sito web per visualizzazione doc. & \glslink{programmatore}{Programmatore} & 1 & 1 & Completata \\
    \hline
    Inizio stesura \glslink{glossario}{Glossario} & \glslink{analista}{Analista} & 3 & 3 & Completata \\
    \hline
    Inizio stesura \glslink{analisi-dei-requisiti-adr}{Analisi dei Requisiti} & \glslink{analista}{Analista} & 6 & 7 & Completata\\
    \hline
    Inizio stesura \glslink{norme-di-progetto}{Norme di Progetto} & \glslink{amministratore}{Amministratore} & 6 & 6 & Completata \\
    \hline
    Costruzione Mock-up & \glslink{progettista}{Progettista} & 4 & 5 & Completata \\
    \hline
    Verifica della documentazione prodotta & \glslink{verificatore}{Verificatore} & 4 & 3 & Completata \\
    \hline
    Redazione verbali interni durante le riunioni & \glslink{responsabile}{Responsabile} & 5 & 5 & Completata \\
    \hline
    \rowcolor{gray!20}
    \textbf{Totale} & & \textbf{29} & \textbf{30} & \\
    \hline
    \end{tabular}%
    }
    \caption{Dettaglio delle attività, ore stimate vs effettive e stato di completamento - Sprint 1}
    \label{tab:task_sprint1}
\end{table}

\paragraph{Considerazioni (Sprint Review e Retrospective)}
\begin{itemize}
    \item \textbf{Analisi dei Risultati:} Avvio efficace delle attività fondanti del progetto, con la configurazione dell'ambiente di lavoro, la prima stesura della documentazione principale (\glslink{glossario}{Glossario}, \glslink{analisi-dei-requisiti-adr}{Analisi dei Requisiti}, \glslink{norme-di-progetto}{Norme di Progetto}) e la realizzazione della prima versione del mock-up grafico.
    \item \textbf{Azioni Correttive:} Maggiore accortezza nelle stime per le attività di Analisi e Progettazione durante i prossimi \glslink{sprint}{Sprint} Planning; introduzione dell'obbligo di avviso tempestivo al \glslink{responsabile}{Responsabile} in caso di ritardi per consentire una rapida redistribuzione del carico.
\end{itemize}

\paragraph{Considerazioni (Sprint Retrospective)}
\textbf{Problematiche Riscontrate:} Lieve rallentamento fisiologico per i ruoli di \glslink{analista}{Analista} e \glslink{progettista}{Progettista}, dovuto alla curva di apprendimento nell'utilizzo dei tool grafici e alla fase di rodaggio nella stesura dei requisiti.\\

I rischi riscontrati sono stati i seguenti:
\begin{itemize}
    \item \textbf{RT1:} La progettazione del mock-up, ha richiesto all'inizio uno studio su nuove tecnologie. Il problema è stato sviato decidendo di lavorare con tecnologie e linguaggi conosciuti da tutti in modo da non spendere tempo a imparare nuove tecnologie in questa fase di apprendimento, e concentrarsi maggiormente sulla stesura dei documenti. 
\end{itemize}
    \newpage

\subsubsubsection{Secondo Periodo}

\label{sec:rtb_secondoperiodo}
\begin{table}[H]
    \centering
    \renewcommand{\arraystretch}{1.5} % Aumenta l'altezza delle righe
    \begin{tabularx}{\textwidth}{|X|X|}
    \hline \rowcolor{Dandelion!50}
    \textbf{Data Inizio} & \textbf{Data Fine}\\
    \hline
    13/04/2026 & 26/04/2026\\
    \hline

    \end{tabularx}
\end{table}

Durante il secondo \glslink{sprint}{Sprint}, l'obiettivo principale del team è stato il consolidamento e l'allineamento della documentazione per prevenire future incomprensioni sui requisiti. Si è proceduto con l'avvio della stesura del \textit{\glslink{piano-di-progetto}{Piano di Progetto}} e con il continuo aggiornamento dell'\textit{\glslink{analisi-dei-requisiti-adr}{Analisi dei Requisiti}} (tramite l'estensione dei casi d'uso) e del \textit{\glslink{glossario}{Glossario}}. Sul fronte progettuale, il mock-up è stato ultimato e presentato all'azienda \glslink{committente-proponente}{proponente}. A seguito del feedback ricevuto, sono state recepite le richieste per alcune modifiche minori; la revisione finale e la relativa approvazione del mock-up sono state riprogrammate per lo \glslink{sprint}{Sprint} successivo.\\

\begin{center}
\noindent\textbf{Attivitá Pianificate \textit{(\glslink{sprint}{Sprint} Backlog)}}\\
\end{center}

\begin{table}[H]
    \centering
    \renewcommand{\arraystretch}{1.5}
    \resizebox{\textwidth}{!}{%
    \begin{tabular}{|l|c|c|c|c|}
    \hline
    \rowcolor{Dandelion!50}
    \textbf{Descrizione Incarico} & \textbf{Ruolo Assegnato} & \textbf{Ore Previste} & \textbf{Ore Effettive} & \textbf{Stato} \\
    \hline
    Inizio stesura \glslink{piano-di-progetto}{Piano di Progetto} & \glslink{responsabile}{Responsabile} & 3 & 3 & Completata \\
    \hline
    Aggiornamento \glslink{glossario}{Glossario} & \glslink{analista}{Analista} & 3 & 4 & Completata \\
    \hline
    Aggiornamento \glslink{norme-di-progetto}{Norme di Progetto} & \glslink{amministratore}{Amministratore} & 6 & 5 & Completata \\
    \hline
    Aggiornamento \glslink{analisi-dei-requisiti-adr}{Analisi dei Requisiti} & \glslink{analista}{Analista} & 5 & 6 & In corso \\
    \hline
    Rifinitura Mock-up & \glslink{progettista}{Progettista} & 2 & 2 & Completata \\
    \hline
    Verifica della documentazione prodotta & \glslink{verificatore}{Verificatore} & 3 & 5 & Completata \\
    \hline
    Redazione verbali interni e \glslink{verbale-esterno}{verbale esterno} & \glslink{responsabile}{Responsabile} & 2 & 2 & Completata \\
    \hline
    \rowcolor{gray!20}
    \textbf{Totale} & & \textbf{24} & \textbf{27} & \\
    \hline
    \end{tabular}%
    }
    \caption{Dettaglio delle attività, ore stimate vs effettive e stato di completamento - Sprint 2}
    \label{tab:task_sprint2}
\end{table}


\paragraph{Considerazioni (Sprint Review)}
\begin{itemize}
    \item \textbf{Analisi dei Risultati:} Ultimazione e presentazione del mock-up grafico all'azienda \glslink{committente-proponente}{proponente}; buon avanzamento nel consolidamento e nell'allineamento della documentazione per prevenire ambiguità.
    \item \textbf{Azioni Correttive:} Ottimizzazione della \glslink{pianificazione}{pianificazione} interna per riassorbire lo sforzo orario extra e impegno del team nell'affinare le metriche di stima a partire dall'iterazione successiva.
\end{itemize}

\paragraph{Considerazioni (Sprint Retrospective)}

\textbf{Problematiche Riscontrate:} Emersione di richieste per modifiche minori sul mock-up da parte del \glslink{committente-proponente}{proponente} e incremento imprevisto del carico di lavoro per i ruoli di \glslink{analista}{Analista} (estensione casi d'uso) e \glslink{verificatore}{Verificatore}. Le attività di \glslink{analisi-dei-requisiti-adr}{Analisi dei Requisiti} hanno subito un momentaneo stallo a causa della necessità di acquisire dal \glslink{committente-proponente}{proponente} dei file di esempio (es. cedolini, comunicazioni). Tali documenti si sono rivelati un prerequisito indispensabile per poter estrarre e mappare correttamente i dati necessari alla descrizione dettagliata dei casi d'uso, generando un rallentamento rispetto alle stime iniziali.
Come azione correttiva, il team ha riassegnato temporaneamente gli Analisti ad altre attività documentali (ovvero la correzione dello stesso documento) in attesa della ricezione dei file, per poi concentrare lo sforzo sui casi d'uso non appena i materiali sono stati resi disponibili.\\

I rischi riscontrati sono stati i seguenti:
\begin{itemize}
    \item \textbf{RR1:} Per quanto riguarda l'estenzione dei casi d'uso. Siccome, nel seguente caso, il rallentamento non è stato eccessivo, i nuovi casi d'uso sono stati aggiunti senza comportare un ritardo nelle scadenze. Tuttavia, è importante sottolineare che, qualora si verificassero situazioni analoghe in futuro, potrebbe essere necessario rivedere le scadenze per garantire un adeguato tempo di sviluppo e mantenere la qualità del lavoro.
\end{itemize}
\newpage

\subsubsubsection{Terzo Periodo}
\label{sec:rtb_terzoperiodo}
\begin{table}[H]
    \centering
    \renewcommand{\arraystretch}{1.5} % Aumenta l'altezza delle righe
    \begin{tabularx}{\textwidth}{|X|X|}
    \hline \rowcolor{Dandelion!50}
    \textbf{Data Inizio} & \textbf{Data Fine}\\
    \hline
    27/04/2026 & 10/05/2026\\
    \hline

    \end{tabularx}
\end{table}

Durante il terzo \glslink{sprint}{Sprint}, il team ha proseguito attivamente con la stesura della documentazione, concentrandosi in particolar modo sul consolidamento dell'\textit{\glslink{analisi-dei-requisiti-adr}{Analisi dei Requisiti}} (estensione dei casi d'uso e diagrammazione), sull'ampliamento del \textit{\glslink{piano-di-progetto}{Piano di Progetto}} e sull'aggiornamento continuo del \textit{\glslink{glossario}{Glossario}}. Inoltre, è stata iniziata la stesura del \textit{\glslink{piano-di-qualifica}{Piano di Qualifica}}. Dal punto di vista tecnico, è stata ultimata la revisione del mock-up per il \glslink{committente-proponente}{proponente} ed è iniziata la fase di autoformazione sui framework scelti per lo sviluppo (\glslink{laravel}{Laravel} e \glslink{filament}{Filament}).

\begin{center}
\noindent\textbf{Attivitá Pianificate \textit{(\glslink{sprint}{Sprint} Backlog)}}\\
\end{center}

\begin{table}[H]
    \centering
    \renewcommand{\arraystretch}{1.5}
    \resizebox{\textwidth}{!}{%
    \begin{tabular}{|l|c|c|c|c|}
    \hline
    \rowcolor{Dandelion!50}
    \textbf{Descrizione Incarico} & \textbf{Ruolo Assegnato} & \textbf{Ore Previste} & \textbf{Ore Effettive} & \textbf{Stato} \\
    \hline
    Aggiornamento \glslink{piano-di-progetto}{Piano di Progetto} & \glslink{responsabile}{Responsabile} & 1 & 1 & Completata \\
    \hline
    Aggiornamento \glslink{glossario}{Glossario} & \glslink{analista}{Analista} & 2 & 3 & Completata \\
    \hline
    Aggiornamento \glslink{norme-di-progetto}{Norme di Progetto} & \glslink{amministratore}{Amministratore} & 3 & 3 & Completata \\
    \hline
    Aggiornamento \glslink{analisi-dei-requisiti-adr}{Analisi dei Requisiti} & \glslink{analista}{Analista} & 7 & 9 & Completata \\
    \hline
    Inizio stesura \glslink{piano-di-qualifica}{Piano di Qualifica} & \glslink{responsabile}{Responsabile} & 1 & 1 & Completata \\
    \hline
    Rifinitura Mock-up & \glslink{programmatore}{Programmatore} & 0 & 2 & Completata \\
    \hline
    Studio delle tecnologie adottate & \glslink{progettista}{Progettista} & 12 & 11 & Completata \\
    \hline
    Verifica della documentazione prodotta & \glslink{verificatore}{Verificatore} & 3 & 3 & Completata \\
    \hline
    Redazione verbali interni e \glslink{verbale-esterno}{verbale esterno} & \glslink{responsabile}{Responsabile} & 2 & 2 & Completata \\
    \hline
    \rowcolor{gray!20}
    \textbf{Totale} & & \textbf{31} & \textbf{35} & \\
    \hline
    \end{tabular}%
    }
    \caption{Dettaglio delle attività, ore stimate vs effettive e stato di completamento - Sprint 3}
    \label{tab:task_sprint3}
\end{table}

\paragraph{Considerazioni (Sprint Review)}
\begin{itemize}
    \item \textbf{Analisi dei Risultati:} Approvazione finale del mock-up da parte del \glslink{committente-proponente}{proponente}, forte progressione nella stesura dei casi d'uso (\glslink{analisi-dei-requisiti-adr}{AdR}) e avvio efficace della fase di autoformazione sui framework tecnologici \glslink{laravel}{Laravel} e \glslink{filament}{Filament}.
    \item \textbf{Azioni Correttive:} Gestione flessibile e tempestiva del budget orario, attuando una riduzione mirata sulle ore del \glslink{progettista}{Progettista} per supportare i ruoli che hanno portato via più tempo del necessario (\glslink{analista}{Analista} e \glslink{programmatore}{Programmatore}).
\end{itemize}

\paragraph{Considerazioni (Sprint Retrospective)}
\textbf{Problematiche Riscontrate:} Sottostima iniziale dell'intensità del lavoro richiesto per la diagrammazione dei requisiti ed emersione di attività tecniche non preventivate (ruolo del \glslink{programmatore}{Programmatore}) necessarie alla rifinitura del mock-up. 
\newpage

\subsubsubsection{Quarto Periodo}

\label{sec:rtb_quartoperiodo}
\begin{table}[H]
    \centering
    \renewcommand{\arraystretch}{1.5} 
    \begin{tabularx}{\textwidth}{|X|X|}
    \hline \rowcolor{Dandelion!50}
    \textbf{Data Inizio} & \textbf{Data Fine}\\
    \hline
    11/05/2026 & 24/05/2026\\
    \hline

    \end{tabularx}
\end{table}

Il quarto periodo di \glslink{sprint}{Sprint} ha rappresentato il momento di convergenza e finalizzazione di tutte le attività legate alla fase di \glslink{rtb-requirements-and-technology-baseline}{RTB}. Sul piano documentale, il team si è focalizzato sulla chiusura formale dell'\glslink{analisi-dei-requisiti-adr}{Analisi dei Requisiti} e del \glslink{glossario}{Glossario}, integrando gli ultimi raffinamenti emersi dalle verifiche interne. Il nucleo operativo dell'iterazione è stato tuttavia assorbito dalle attività tecniche, con la definizione teorica e la successiva implementazione pratica del PoC. Sfruttando lo stack tecnologico concordato il gruppo ha sviluppato un prototipo funzionante mirato a validare la fattibilità delle architetture e a minimizzare i rischi tecnologici individuati. 

A conclusione dell'iterazione, il team ha organizzato un incontro ufficiale con l'azienda \glslink{committente-proponente}{proponente} Eggon per presentare una demo operativa del PoC. Il colloquio ha permesso di mostrare le soluzioni software realizzate, ricevendo un riscontro estremamente positivo da parte del \glslink{committente-proponente}{committente}; questo passaggio ha sancito il superamento dei criteri di accettazione interni e ha confermato la solidità del lavoro svolto in preparazione al colloquio di revisione \glslink{rtb-requirements-and-technology-baseline}{RTB}.

\begin{center}
\noindent\textbf{Attivitá Pianificate \textit{(\glslink{sprint}{Sprint} Backlog)}}\\
\end{center}
\begin{table}[H]
    \centering
    \renewcommand{\arraystretch}{1.5}
    \resizebox{\textwidth}{!}{%
    \begin{tabular}{|l|c|c|c|c|}
    \hline
    \rowcolor{Dandelion!50}
    \textbf{Descrizione Incarico} & \textbf{Ruolo Assegnato} & \textbf{Ore Previste} & \textbf{Ore Effettive} & \textbf{Stato} \\
    \hline
    Aggiornamento \glslink{piano-di-progetto}{Piano di Progetto} & \glslink{responsabile}{Responsabile} & 1 & 1 & Completata \\
    \hline
    Stesura e incremento \glslink{piano-di-qualifica}{Piano di Qualifica} & \glslink{responsabile}{Responsabile} & 2 & 2 & Completata \\
    \hline
    Aggiornamento \glslink{norme-di-progetto}{Norme di Progetto} & \glslink{amministratore}{Amministratore} & 1 & 1 & Completata \\
    \hline
    Conclusione \glslink{analisi-dei-requisiti-adr}{Analisi dei Requisiti} e \glslink{glossario}{Glossario} & \glslink{analista}{Analista} & 1 & 1 & Completata \\
    \hline
    Definizione del PoC & \glslink{progettista}{Progettista} & 13 & 15 & Completata \\
    \hline
    Implementazione del PoC & \glslink{programmatore}{Programmatore} & 16 & 18 & Completata \\
    \hline
    Verifica della documentazione e del PoC & \glslink{verificatore}{Verificatore} & 8 & 11 & Completata \\
    \hline
    Redazione verbali interni ed esterni & \glslink{responsabile}{Responsabile} & 1 & 1 & Completata \\
    \hline
    \rowcolor{gray!20}
    \textbf{Totale} & & \textbf{43} & \textbf{50} & \\
    \hline
    \end{tabular}%
    }
    \caption{Dettaglio delle attività, ore stimate vs effettive e stato di completamento - Sprint 4}
    \label{tab:task_sprint4}
\end{table}

\paragraph{Considerazioni (Sprint Review)}
\begin{itemize}
    \item \textbf{Analisi dei Risultati:} Chiusura definitiva dell'\glslink{analisi-dei-requisiti-adr}{Analisi dei Requisiti} e sviluppo completato del Proof of Concept (PoC) utilizzando le tecnologie designate; Completati \glslink{norme-di-progetto}{Norme di progetto}, \glslink{piano-di-qualifica}{piano di qualifica}, \glslink{piano-di-progetto}{piano di progetto} e \glslink{glossario}{glossario} in vista dell'\glslink{rtb-requirements-and-technology-baseline}{RTB}.
    \item \textbf{Azioni Correttive:} Trattandosi dell'ultimo \glslink{sprint}{Sprint} della fase \glslink{rtb-requirements-and-technology-baseline}{RTB}, il gruppo utilizzerà le criticità per il \glslink{preventivo}{preventivo} delle ore del PoC per stilare una \glslink{pianificazione}{pianificazione} oraria post-\glslink{rtb-requirements-and-technology-baseline}{RTB} estremamente accurata, prevedendo \textit{buffer} di \glslink{rischio}{rischio} più ampi per la fase di coding.
\end{itemize}

\paragraph{Considerazioni (Sprint Retrospective)}
\textbf{Problematiche Riscontrate:} Sul piano tecnico, lo sviluppo del PoC ha subito un momentaneo stallo a causa di una dipendenza esterna legata all'infrastruttura cloud di \textbf{AWS}. Il team si è trovato bloccato in attesa di poter richiedere e ottenere le credenziali e le risorse necessarie direttamente dall'azienda \glslink{committente-proponente}{proponente}. Questo periodo di attesa forzata ha costretto i ruoli tecnici (\glslink{progettista}{Progettista}, \glslink{programmatore}{Programmatore} e \glslink{verificatore}{Verificatore}) a condensare l'intera fase di implementazione e collaudo in un arco di tempo molto più ristretto rispetto a quanto pianificato. Per recuperare il tempo latente e garantire la consegna del PoC entro i termini dello \glslink{sprint}{Sprint}, il team ha dovuto svolgere sessioni di lavoro straordinarie e intensive, generando così lo scostamento complessivo di +7 ore.\\
I rischi riscontrati sono stati i seguenti:
\begin{itemize}
\item \textbf{RT1:} Per quanto riguarda lo studio delle nuove tecnologie per lo sviluppo del PoC. Il problema è stato mitigato, come già descritto inizialmente, dedicando del tempo per imparare le nuove tcnologie, cercando il più possibile una documentazione chiara e sufficientemente dettagliata. 
\end{itemize}

\subsection{Pianificazione a Lungo Termine (Verso la PB)}
La programmazione analitica dei singoli \glslink{sprint}{Sprint} e delle scadenze intermedie verrà ufficializzata e integrata nella release v2.0.0 del presente documento, dopo aver superato la \glslink{rtb-requirements-and-technology-baseline}{RTB}. Nonostante ciò, il team ha delineato subito una strategia di sviluppo a lungo termine da poter seguire. Questo approccio, articolato in 4 macro-fasi sequenziali, è concepito per assicurare un rilascio incrementale, isolando le criticità architetturali.

\subsubsection{Macro-fase 1: Consolidamento Architetturale e Infrastrutturale (Post-RTB)}
Questa fase segna il passaggio dallo sviluppo sperimentale PoC alla progettazione ingegneristica rigorosa del sistema.
\begin{itemize}
    \item \textbf{Obiettivo:} Stabilizzare l'architettura di base e predisporre l'ambiente cloud definitivo, garantendo che le fondamenta del software siano scalabili e sicure per il trattamento di dati aziendali.
    \item \textbf{Attività chiave:} Stesura del \textit{Technical Design} definitivo (diagrammi UML avanzati), configurazione strutturata dei database e dello storage cloud per i documenti, implementazione dei protocolli di sicurezza per i dati sensibili in conformità al \glslink{gdpr}{GDPR} e setup della pipeline automatizzata di CI/CD.
\end{itemize}

\subsubsection{Macro-fase 2: Sviluppo del Core del Prodotto (Fase MVP)}
Fase interamente focalizzata sulla realizzazione del nucleo funzionale essenziale per la piattaforma, al fine di soddisfare tutti i requisiti classificati come obbligatori.
\begin{itemize}
    \item \textbf{Obiettivo:} Rilasciare un \textit{\glslink{mvp-minimum-viable-product}{Minimum Viable Product}} (\glslink{mvp-minimum-viable-product}{MVP}) stabile, in grado di dimostrare l'interoperabilità di base tra le interfacce utente e i servizi di backend.
    \item \textbf{Attività chiave:} Sviluppo delle API principali di comunicazione, implementazione delle componenti core per la \glslink{dashboard}{dashboard} amministrativa e le interfacce lato dipendente, e prima integrazione stabile dei motori di \glslink{intelligenza-artificiale-ai}{Intelligenza Artificiale} per l'elaborazione dei testi e il riconoscimento documentale di base.
\end{itemize}

\subsubsection{Macro-fase 3: Estensione Funzionale, Ottimizzazione e Sicurezza}
Fase orientata al completamento dei requisiti desiderabili e opzionali, focalizzandosi sull'efficienza e sulla rifinitura dell'esperienza utente (UX).
\begin{itemize}
    \item \textbf{Obiettivo:} Raggiungere la piena copertura funzionale del \glslink{capitolato}{capitolato}, ottimizzando i costi operativi legati alle chiamate \glslink{intelligenza-artificiale-ai}{AI} e perfezionando le logiche di automazione.
    \item \textbf{Attività chiave:} Estensione delle funzionalità \glslink{intelligenza-artificiale-ai}{AI} (raffinamento dei \glslink{prompt}{prompt} e dei meccanismi di smistamento automatico), implementazione di sistemi di caching per minimizzare le richieste esterne, definizione granulare del controllo accessi basato sui ruoli (\glslink{rbac-role-based-access-control}{RBAC}) e rifinitura grafica delle interfacce \textit{\glslink{nexum}{Nexum}}.
\end{itemize}

\subsubsection{Macro-fase 4: Validazione, Quality Assurance e Rilascio}
L'ultimo periodo prima della consegna ufficiale, dedicato esclusivamente alle attività di verifica approfondita e alla preparazione del deployment finale.
\begin{itemize}
    \item \textbf{Obiettivo:} Certificare la conformità totale del prodotto rispetto \glslink{intelligenza-artificiale-ai}{ai} requisiti qualitativi, prestazionali e di sicurezza concordati con la \glslink{committente-proponente}{proponente}.
    \item \textbf{Attività chiave:} Esecuzione intensiva dei \textit{System Test} e dei \glslink{test-di-integrazione}{test di integrazione}, sessioni di \textit{User Acceptance Test} (UAT) in coordinamento con Eggon, stesura e validazione della manualistica finale (Manuale Utente e Manuale \glslink{amministratore}{Amministratore}) e rilascio del codice in ambiente di produzione.
\end{itemize}

\newpage
\section{Preventivo e Consuntivo dei Periodi}
\label{sec:preventivo_e_consuntivo}
\subsection{Primo Periodo: Sprint 30/03/2026 - 12/04/2026}
\label{sec:pec_primoperiodo}

\vspace{0.5em}

\begin{table}[H]
\centering
\label{tab:preventivo-sprint1}
\vspace{2mm}

\begin{center}
    
\noindent\textbf{\glslink{preventivo}{Preventivo}}
\end{center}

\begin{tabular}{|l|c|c|c|c|c|c|>{\cellcolor{Apricot!50}}c|}
\hline
\textbf{Componente del Team} & \textbf{Resp.} & \textbf{Amm.} & \textbf{Ana.} & \textbf{Proget.} & \textbf{Progr.} & \textbf{Ver.} & \textbf{Totale} \\
\hline
Angelica Gastaldello & 5 & - & 1 & - & - & - & 6\\
Eleonora Bellet & - & 3 & - & - & - & - & 3\\
Alex Oloni & - & - & 1 & 2 & - & - & 3\\
Diego Belluz & - & 3 & - & - & - & - & 3\\
Angela Maule & - & - & 2 & - & - & 1 & 3\\
Leonardo Soligo & - & - & 3 & 2 & 1 & - & 6\\
V.Y. Kaneda Fernandes & - & - & 2 & - & - & 3 & 5\\
\hline
\hiderowcolors
\rowcolor{green!50!black!20!white}
\textbf{Totale} & 5 & 6 & 9 & 4 & 1 & 4 & 29\\
\hline
\end{tabular}
\caption{Preventivo ore per componente - Primo Sprint}
\end{table}

\begin{center}
    
\noindent\textbf{\glslink{consuntivo}{Consuntivo}}
\end{center}

\begin{table}[H]
\centering
\vspace{2mm}
\begin{tabular}{|l|c|c|c|c|c|c|>{\cellcolor{Apricot!50}}c|}
\hline
\textbf{Componente del Team} & \textbf{Resp.} & \textbf{Amm.} & \textbf{Ana.} & \textbf{Proget.} & \textbf{Progr.} & \textbf{Ver.} & \textbf{Totale} \\
\hline
Angelica Gastaldello & 5 & - & 2 & - & - & - & 7\\
Eleonora Bellet & - & 3 & - & - & - & - & 3\\
Alex Oloni & - & - & 1 & 2 & - & - & 3\\
Diego Belluz & - & 3 & - & - & - & - & 3\\
Angela Maule & - & - & 3 & - & - & 1 & 4\\
Leonardo Soligo & - & - & 2 & 3 & 1 & - & 6 \\
V.Y. Kaneda Fernandes & - & - & 2 & - & - & 2 & 4\\
\hline
\hiderowcolors
\rowcolor{green!50!black!20!white}
\textbf{Totale} & 5 & 6 & 10 & 5 & 1 & 3 & 30\\
\hline
\end{tabular}
\caption{Consuntivo ore per componente - Primo Sprint}
\end{table}

\begin{figure}[h]
    \centering
    \begin{tikzpicture}
        \pie[
            text=legend,
            radius=3,
            color={red!60, blue!60, yellow!60, green!60, purple!60, orange!60}
        ]{
            7.7/\glslink{responsabile}{Responsabile},
            23/\glslink{amministratore}{Amministratore},
            38.5/\glslink{analista}{Analista},
            19.2/\glslink{progettista}{Progettista},
            3.9/\glslink{programmatore}{Programmatore},
            7.7/\glslink{verificatore}{Verificatore}
        }
    \end{tikzpicture}
    \caption{Grafico Consuntivo Costi - Primo Sprint}
\end{figure}

\vspace{0.5cm}
\begin{center}
    
\noindent\textbf{Bilancio Economico}
\end{center}

\begin{table}[H]
\centering
\label{tab:risorse-rimanenti}
\vspace{2mm}
\begin{tabular}{|l|c|c|c|c|c|c|>{\cellcolor{Apricot!50}}c|}
\hline
\textbf{Costo} & \textbf{Resp.} & \textbf{Amm.} & \textbf{Ana.} & \textbf{Proget.} & \textbf{Progr.} & \textbf{Ver.} & \textbf{Totale} \\
\hline
Costo Orario &  30 \euro &  20 \euro &  25 \euro&  25 \euro &  15 \euro &  15 \euro & /\\
Costo Totale &  150 \euro& 120 \euro & 250 \euro & 125 \euro & 15 \euro & 45 \euro & 705 \euro\\
\hline
\hiderowcolors
\rowcolor{green!50!black!20!white}
\textbf{Saldo fine periodo} & 1530 \euro & 900 \euro & 1250 \euro & 3450 \euro & 2385 \euro & 2355 \euro & \cellcolor{yellow!50} 11.870 \euro  \\
\hline
\end{tabular}
\caption{Riepilogo costi e budget rimanente - Primo Sprint }
\end{table}

\subsubsection{Aggiornamento del Preventivo a Finire - Sprint 1}
Al termine del primo \glslink{sprint}{Sprint}, viene effettuato il calcolo del \textit{\glslink{preventivo}{Preventivo} a Finire} (EAC - \textit{Estimate at Completion}) per monitorare proattivamente l'andamento economico del progetto. Questa metrica proietta il costo totale finale sulla base degli scostamenti registrati nell'iterazione appena conclusa.

Rispetto al \glslink{preventivo}{preventivo} iniziale dello \glslink{sprint}{Sprint} 1, si è registrato un lieve scostamento in eccesso per i ruoli di \glslink{analista}{Analista} (+1 ora) e \glslink{progettista}{Progettista} (+1 ora), parzialmente compensato da un risparmio nel ruolo di \glslink{verificatore}{Verificatore} (-1 ora). Di seguito viene presentata la proiezione aggiornata sull'intero ciclo di vita del progetto.

\begin{table}[H]
    \centering
    \renewcommand{\arraystretch}{1.5}
    \resizebox{\textwidth}{!}{
    \begin{tabular}{|l|c|c|c|c|c|}
    \hline
    \textbf{Ruolo} & \textbf{Ore Prev. Iniziali} & \textbf{Scostamento Ore} & \textbf{EAC (Ore)} & \textbf{EAC (Costo)} & \textbf{Varianza} \\
    \hline
    \glslink{responsabile}{Responsabile} & 56 & +0 & 56 & 1.680,00 \euro & 0,00 \euro \\
    \hline
    \glslink{amministratore}{Amministratore} & 52 & +0 & 52 & 1.020,00 \euro & 0,00 \euro\\
    \hline
    \glslink{analista}{Analista} & 59 & +1 & 60 & 1.500,00 \euro & -25,00 \euro\\
    \hline
    \glslink{progettista}{Progettista} & 143 & +1 & 144 & 3.575,00 \euro & -25,00 \euro\\
    \hline
    \glslink{programmatore}{Programmatore} & 160 & +0 & 160 & 2.400,00 \euro & 0,00 \euro\\
    \hline
    \glslink{verificatore}{Verificatore} & 160 & -1 & 159 & 2.400,00 \euro & +15,00 \euro \\
    \hline
    \rowcolor{green!50!black!20!white}
   \textbf{Totale} & \textbf{630} & \textbf{+1} & \textbf{631} & \textbf{12.605,00 \euro} & \textbf{-35,00 \euro} \\
    \hline
    \end{tabular}%
    }
    \caption{Aggiornamento del Preventivo a Finire e Varianza al termine dello Sprint 1}
    \label{tab:eac_sprint1}
\end{table}

Come si nota dalla Tabella~\ref{tab:eac_sprint1}, il bilancio a finire proietta un totale di 631 ore con un costo complessivo stimato di € 12.610,00. La Varianza a Finire (VAC) risulta essere leggermente negativa e pari a € -35,00. Essendo il primo \glslink{sprint}{Sprint} di progetto, questo scostamento è considerato del tutto fisiologico e legato alla fase di rodaggio del team; esso verrà tenuto in considerazione nelle prossime pianificazioni per riportare il bilancio in perfetto equilibrio.

\newpage
\subsection{Secondo Periodo: Sprint 13/04/2026 - 26/04/2026}
\label{sec:pec_secondoperiodo}

\vspace{0.5em}

\begin{table}[H]
\centering
\vspace{2mm}

\begin{center}
    
\noindent\textbf{\glslink{preventivo}{Preventivo}}
\end{center}

\begin{table}[H]
\centering
\vspace{2mm}
\begin{tabular}{|l|c|c|c|c|c|c|>{\cellcolor{Apricot!50}}c|}
\hline
\textbf{Componente del Team} & \textbf{Resp.} & \textbf{Amm.} & \textbf{Ana.} & \textbf{Proget.} & \textbf{Progr.} & \textbf{Ver.} & \textbf{Totale} \\
\hline
Angelica Gastaldello & - & - & 3 & - & - & - & 3\\
Eleonora Bellet & - & 3 & - & - & - & 1 & 4\\
Alex Oloni & - & 3 & - & - & - & - & 3\\
Diego Belluz & - & - & 2 & - & - & - & 2\\
Angela Maule & 5 & - & 1 & - & - & - & 6\\
Leonardo Soligo & - & - & - & 2 & - & - & 2\\
V.Y. Kaneda Fernandes & - & - & 2 & - & - & 2 & 4\\
\hline
\hiderowcolors
\rowcolor{green!50!black!20!white}
\textbf{Totale} & 5 & 6 & 8 & 2 & 0 & 3 & 24\\
\hline
\end{tabular}
\caption{Preventivo ore per componente - Secondo Sprint}
\end{table}

\begin{center}
    
\noindent\textbf{\glslink{consuntivo}{Consuntivo}}
\end{center}

\begin{tabular}{|l|c|c|c|c|c|c|>{\cellcolor{Apricot!50}}c|}
\hline
\textbf{Componente del Team} & \textbf{Resp.} & \textbf{Amm.} & \textbf{Ana.} & \textbf{Proget.} & \textbf{Progr.} & \textbf{Ver.} & \textbf{Totale} \\
\hline
Angelica Gastaldello & - & - & 4 & - & - & - & 4\\
Eleonora Bellet & - & 3 & - & - & - & 2 & 5\\
Alex Oloni & - & 2 & - & - & - & - & 2\\
Diego Belluz & - & - & 2 & - & - & - & 2\\
Angela Maule & 5 & - & 2 & - & - & - & 7\\
Leonardo Soligo & - & - & - & 2 & - & - & 2\\
V.Y. Kaneda Fernandes & - & - & 2 & - & - & 3 & 5\\
\hline
\hiderowcolors
\rowcolor{green!50!black!20!white}
\textbf{Totale} & 5 & 5 & 10 & 2 & 0 & 5 & 27\\
\hline
\end{tabular}
\caption{Consuntivo ore per componente - Secondo Sprint}
\end{table}

\begin{figure}[h]
    \centering
    \begin{tikzpicture}
        \pie[
            text=legend,
            radius=3,
            color={red!60, blue!60, yellow!60, green!60, orange!60}
        ]{
            21.7/\glslink{responsabile}{Responsabile},
            21.7/\glslink{amministratore}{Amministratore},
            43.5/\glslink{analista}{Analista},
            8.7/\glslink{progettista}{Progettista},
            4.4/\glslink{verificatore}{Verificatore}
        }
    \end{tikzpicture}
    \caption{Grafico Consuntivo Costi - Secondo Sprint}
\end{figure}


\vspace{0.5cm}
\begin{center}
    
\noindent\textbf{Bilancio Economico}
\end{center}

\begin{table}[H]
\centering
\vspace{2mm}
\begin{tabular}{|l|c|c|c|c|c|c|>{\cellcolor{Apricot!50}}c|}
\hline
\textbf{Costo} & \textbf{Resp.} & \textbf{Amm.} & \textbf{Ana.} & \textbf{Proget.} & \textbf{Progr.} & \textbf{Ver.} & \textbf{Totale} \\
\hline
Costo Orario &  30 \euro &  20 \euro &  25 \euro&  25 \euro &  15 \euro &  15 \euro & /\\
Costo Totale &  150 \euro& 100 \euro & 250 \euro & 50 \euro & 0 \euro & 75 \euro & 625\euro\\
\hline
\hiderowcolors
\rowcolor{green!50!black!20!white}
\textbf{Saldo fine periodo} & 1380 \euro & 800 \euro & 1000 \euro & 3400 \euro & 2385 \euro & 2280 \euro & \cellcolor{yellow!50} 11.245 \euro  \\
\hline
\end{tabular}
\caption{Riepilogo costi e budget rimanente - Secondo Sprint}
\end{table}

\subsubsection{Aggiornamento del Preventivo a Finire - Sprint 2}
Al termine del secondo \glslink{sprint}{Sprint}, il \textit{\glslink{preventivo}{Preventivo} a Finire} (EAC) è stato ricalcolato tenendo conto degli scostamenti cumulativi registrati nelle prime due iterazioni di progetto. 

Nello \glslink{sprint}{Sprint} appena concluso si è verificato un lieve scostamento in eccesso rispetto alle stime, principalmente dovuto all'aumento delle ore necessarie per l'\glslink{analisi-dei-requisiti-adr}{Analisi dei Requisiti} e per le attività di Verifica, parzialmente mitigate da un risparmio nel ruolo di \glslink{amministratore}{Amministratore}. Di seguito viene presentata la proiezione economica aggiornata per l'intero progetto.

\begin{table}[H]
    \centering
    \renewcommand{\arraystretch}{1.5}
    \resizebox{\textwidth}{!}{%
    \begin{tabular}{|l|c|c|c|c|c|}
    \hline
    \textbf{Ruolo} & \textbf{Ore Prev. Iniziali} & \textbf{Scostamento Ore} & \textbf{EAC (Ore)} & \textbf{EAC (Costo)} & \textbf{Varianza (VAC)} \\
    \hline
    \glslink{responsabile}{Responsabile} & 56 & +0 & 56 & 1.680,00 \euro & 0,00 \euro\\
    \hline
    \glslink{amministratore}{Amministratore} & 52 & -1 & 51 & 1.020,00 \euro & +20,00 \euro \\
    \hline
    \glslink{analista}{Analista} & 59 & +3 & 62 & 1.550,00 \euro & -75,00 \euro \\
    \hline
    \glslink{progettista}{Progettista} & 143 & +1 & 144 & 3.600,00 \euro & -25,00 \euro\\
    \hline
    \glslink{programmatore}{Programmatore} & 160 & +0 & 160 & 2.400,00 \euro & 0,00 \euro \\
    \hline
    \glslink{verificatore}{Verificatore} & 160 & +1 & 161 & 2.415,00 \euro & -15,00 \euro\\
    \hline
    \rowcolor{green!50!black!20!white}
    \textbf{Totale} & \textbf{630} & \textbf{+4} & \textbf{634} & \textbf{12.665,00 \euro} & \textbf{-95,00 \euro} \\
    \hline
    \end{tabular}%
    }
    \caption{Aggiornamento del Preventivo a Finire e Varianza cumulativa al termine dello Sprint 2}
    \label{tab:eac_sprint2}
\end{table}

Come mostrato nella Tabella~\ref{tab:eac_sprint2}, lo scostamento cumulativo totale si attesta a +4 ore rispetto al piano originale. La Varianza a Finire (VAC) registra un valore di € -95,00. Il team ha preso atto di questa tendenza al rialzo, dovuta principalmente al consolidamento della documentazione; verranno pertanto pianificate delle azioni di ottimizzazione nelle successive iterazioni tecniche per riassorbire tale scostamento e rientrare nel budget stabilito.

\newpage

\subsection{Terzo Periodo: Sprint 27/04/2026 - 10/05/2026}
\label{sec:pec_terzoperiodo}

\vspace{0.5em}

\begin{center}
\noindent\textbf{\glslink{preventivo}{Preventivo}}
\end{center}

\begin{table}[H]
\centering
\vspace{2mm}
\begin{tabular}{|l|c|c|c|c|c|c|>{\cellcolor{Apricot!50}}c|}
\hline
\textbf{Componente del Team} & \textbf{Resp.} & \textbf{Amm.} & \textbf{Ana.} & \textbf{Proget.} & \textbf{Progr.} & \textbf{Ver.} & \textbf{Totale} \\
\hline
Angelica Gastaldello & - & - & 2 & 3 & - & - & 5\\
Eleonora Bellet      & - & - & 1 & 3 & - & 2 & 6\\
Alex Oloni           & 4 & - & - & - & - & - & 4\\
Diego Belluz         & - & 3 & - & 3 & - & - & 6\\
Angela Maule         & - & - & 2 & - & - & - & 2\\
Leonardo Soligo      & - & - & 2 & 3 & - & - & 5\\
V.Y. Kaneda Fernandes & - & - & 2 & - & - & 1 & 3\\
\hline
\hiderowcolors
\rowcolor{green!50!black!20!white}
\textbf{Totale}      & 4 & 3 & 9 & 12 & 0 & 3 & 31\\
\hline
\end{tabular}
\caption{Preventivo ore per componente - Terzo sprint}
\end{table}

\begin{center}
\noindent\textbf{\glslink{consuntivo}{Consuntivo}}
\end{center}

\begin{table}[H]
\centering
\vspace{2mm}
\begin{tabular}{|l|c|c|c|c|c|c|>{\cellcolor{Apricot!50}}c|}
\hline
\textbf{Componente del Team} & \textbf{Resp.} & \textbf{Amm.} & \textbf{Ana.} & \textbf{Proget.} & \textbf{Progr.} & \textbf{Ver.} & \textbf{Totale} \\
\hline
Angelica Gastaldello & - & - & 3 & 4 & - & - & 7\\
Eleonora Bellet & - & - & 3 & 1 & - & 2 & 6\\
Alex Oloni & 4 & - & - & 1 & - & - & 5\\
Diego Belluz & - & 3 & - & 3 & - & - & 6\\
Angela Maule & - & - & 2 & - & - & - & 2\\
Leonardo Soligo & - &  - & 2 & 2 & 2 & - & 6\\
V.Y. Kaneda Fernandes & - & - & 2 & - & - & 1 & 3\\
\hline
\hiderowcolors
\rowcolor{green!50!black!20!white}
\textbf{Totale} & 4 & 3 & 12 & 11 & 2 & 3 & 35\\
\hline
\end{tabular}
\caption{Consuntivo ore per componente - Terzo Sprint}
\end{table}

\begin{figure}[H]
    \centering
    \begin{tikzpicture}
        \pie[
            text=legend,
            radius=3,
            color={red!60, blue!60, yellow!60, green!60, purple!60, orange!60}
        ]{
            14.4/\glslink{responsabile}{Responsabile},
            7.2/\glslink{amministratore}{Amministratore},
            36.1/\glslink{analista}{Analista},
            33.1/\glslink{progettista}{Progettista},
            3.6/\glslink{programmatore}{Programmatore},
            5.4/\glslink{verificatore}{Verificatore}
        }
    \end{tikzpicture}
    \caption{Grafico Consuntivo Costi - Terzo Sprint}
\end{figure}

\vspace{0.5cm}

\begin{center}
\noindent\textbf{Bilancio Economico}
\end{center}

\begin{table}[H]
\centering
\vspace{2mm}
\begin{tabular}{|l|c|c|c|c|c|c|>{\cellcolor{Apricot!50}}c|}
\hline
\textbf{Costo} & \textbf{Resp.} & \textbf{Amm.} & \textbf{Ana.} & \textbf{Proget.} & \textbf{Progr.} & \textbf{Ver.} & \textbf{Totale} \\
\hline
Costo Orario &  30 \euro &  20 \euro &  25 \euro&  25 \euro &  15 \euro &  15 \euro & /\\
Costo Totale &  120 \euro & 60 \euro & 300 \euro & 275 \euro & 30 \euro & 45 \euro &  830\euro\\
\hline
\hiderowcolors
\rowcolor{green!50!black!20!white}
\textbf{Saldo fine periodo} & 1260 \euro & 740 \euro & 700 \euro & 3125 \euro & 2355 \euro & 2235 \euro & \cellcolor{yellow!50} 10.415 \euro  \\
\hline
\end{tabular}
\caption{Riepilogo costi e budget rimanente - Terzo Sprint}
\end{table}

\subsubsection{Aggiornamento del Preventivo a Finire - Sprint 3}
Al termine del terzo \glslink{sprint}{Sprint}, il \glslink{responsabile}{Responsabile} ha proceduto al ricalcolo del \textit{\glslink{preventivo}{Preventivo} a Finire}. Questa analisi predittiva valuta l'impatto degli scostamenti cumulativi registrati fin dall'inizio del progetto rispetto al budget aggiornato in fase di avvio.

Durante quest'ultima iterazione, il team ha compensato lo scostamento in eccesso registrato precedentemente sul ruolo di \glslink{progettista}{Progettista}, riportando il suo bilancio cumulativo in perfetto equilibrio. Tuttavia, il perdurare di un intenso carico di lavoro sull'\glslink{analisi-dei-requisiti-adr}{Analisi dei Requisiti} e l'avvio delle attività di Programmazione per la chiusura del Mock-up e l'inizio del Poc hanno generato un ulteriore assorbimento di risorse. Di seguito viene presentata la proiezione economica aggiornata per l'intero progetto.

\begin{table}[H]
    \centering
    \renewcommand{\arraystretch}{1.5}
    \resizebox{\textwidth}{!}{%
    \begin{tabular}{|l|c|c|c|c|c|}
    \hline
    \textbf{Ruolo} & \textbf{Ore Prev. Iniziali} & \textbf{Scostamento Ore} & \textbf{EAC (Ore)} & \textbf{EAC (Costo)} & \textbf{Varianza (VAC)} \\
    \hline
    \glslink{responsabile}{Responsabile} & 56 & +0 & 56 & 1.680,00 \euro & 0,00 \euro\\
    \hline
    \glslink{amministratore}{Amministratore} & 52 & -1 & 51 & 1.020,00 \euro & +20,00 \euro \\
    \hline
    \glslink{analista}{Analista} & 59 & +6 & 65 & 1.625,00 \euro & -150,00 \euro \\
    \hline
    \glslink{progettista}{Progettista} & 143 & +0 & 143 & 3.575,00 \euro & 0,00 \euro \\
    \hline
    \glslink{programmatore}{Programmatore} & 160 & +2 & 162 & 2.430,00 \euro & -30,00 \euro\\
    \hline
    \glslink{verificatore}{Verificatore} & 160 & +1 & 161 & 2.415,00 \euro & -15,00 \euro \\
    \hline
    \rowcolor{green!50!black!20!white}
    \textbf{Totale} & \textbf{630} & \textbf{+8} & \textbf{638} & \textbf{12.745,00 \euro} & \textbf{-175,00 \euro} \\
    \hline
    \end{tabular}%
    }
    \caption{Aggiornamento del Preventivo a Finire e Varianza cumulativa al termine dello Sprint 3}
    \label{tab:eac_sprint3}
\end{table}

Come evidenziato dalla Tabella~\ref{tab:eac_sprint3}, il bilancio a finire proietta attualmente un totale di 638 ore, con una Varianza (VAC) complessiva di -175,00 \euro dovuta quasi interamente alle stime al ribasso per il ruolo di \glslink{analista}{Analista} nelle fasi iniziali. Tale scostamento impone al team un'attenta ripianificazione nelle prossime settimane: l'ottimizzazione e il risparmio delle ore di \glslink{analista}{Analista} e \glslink{responsabile}{Responsabile} durante la fase di \glslink{pb-product-baseline}{PB} saranno essenziali per garantire un rientro all'interno del tetto massimo di spesa.

\newpage
\subsection{Quarto Periodo: Sprint 11/05/2026 - 24/05/2026}
\label{sec:pec_quartoperiodo}

\vspace{0.5em}

\begin{center}
\noindent\textbf{\glslink{preventivo}{Preventivo}}
\end{center}

\begin{table}[H]
\centering
\vspace{2mm}
\begin{tabular}{|l|c|c|c|c|c|c|>{\cellcolor{Apricot!50}}c|}
\hline
\textbf{Componente del Team} & \textbf{Resp.} & \textbf{Amm.} & \textbf{Ana.} & \textbf{Proget.} & \textbf{Progr.} & \textbf{Ver.} & \textbf{Totale} \\
\hline
Angelica Gastaldello & - & - & - & 2 & 4 & - & 6\\
Eleonora Bellet      & 4 & - & - & - & - & 2 & 6\\
Alex Oloni           & - & - & - & 2 & 4 & - & 6\\
Diego Belluz         & - & 1 & - & 3 & - & 2 & 6\\
Angela Maule         & - & - & 1 & 2 & 4 & - & 7\\
Leonardo Soligo      & - & - & - & 2 & 4 & - & 6\\
V.Y. Kaneda Fernandes & - & - & - & 2 & - & 4 & 6\\
\hline
\hiderowcolors
\rowcolor{green!50!black!20!white}
\textbf{Totale}      & 4 & 1 & 1 & 13 & 16 & 8 & 43\\
\hline
\end{tabular}
\caption{Preventivo ore per componente - Quarto Sprint}
\end{table}

\begin{center}
\noindent\textbf{\glslink{consuntivo}{Consuntivo}}
\end{center}
\begin{table}[H]
\centering
\vspace{2mm}
\begin{tabular}{|l|c|c|c|c|c|c|>{\cellcolor{Apricot!50}}c|}
\hline
\textbf{Componente del Team} & \textbf{Resp.} & \textbf{Amm.} & \textbf{Ana.} & \textbf{Proget.} & \textbf{Progr.} & \textbf{Ver.} & \textbf{Totale} \\
\hline
    Angelica Gastaldello & & & & 3 & 5 & & 8 \\
    \hline
    Eleonora Bellet & 4 & & & & & 3 & 7 \\
    \hline
    Alex Oloni & & & & 2 & 4 & & 6 \\
    \hline
    Diego Belluz & & 1 & & 2 & & 3 & 6 \\
    \hline
    Angela Maule & & & 1 & 3 & 4 & & 8 \\
    \hline
    Leonardo Soligo & & & & 3 & 5 & & 8 \\
    \hline
    V.Y. Kaneda Fernandes & & & & 2 & & 5 & 7 \\
    \hline
\hiderowcolors
\rowcolor{green!50!black!20!white}
    Totale & 4 & 1 & 1 & 15 & 18 & 11 & 50 \\
    \hline
   \end{tabular}
\caption{Consuntivo ore per componente - Quarto sprint}
\end{table}

\begin{figure}[H]
    \centering
    \begin{tikzpicture}
        \pie[
            hide number,
            text=legend,
            radius=3,
            color={red!60, blue!60, yellow!60, green!60, purple!60, orange!60}
        ]{
            12.31/\glslink{responsabile}{Responsabile},
            2.05/\glslink{amministratore}{Amministratore},
            2.56/\glslink{analista}{Analista},
            38.46/\glslink{progettista}{Progettista},
            27.70/\glslink{programmatore}{Programmatore},
            16.92/\glslink{verificatore}{Verificatore}
        }
        \node at (22:1.6) {12.31\%};  
        \node at (130:1.6) {38.46\%};
        \node at (249:1.6) {27.70\%}; 
        \node at (330:1.6) {16.92\%}; 

        \draw[<-] (48:2.9) -- (48:4.2) node[right] {2.05\%};
        
        \draw[<-] (56:2.9) -- (56:4.6) node[above right] {2.56\%};

    \end{tikzpicture}
    \caption{Grafico Consuntivo Costi - Quarto Sprint}
\end{figure}

\begin{center}
\noindent\textbf{Bilancio Economico}
\end{center}

\begin{table}[H]
\centering
\vspace{2mm}
\begin{tabular}{|l|c|c|c|c|c|c|>{\cellcolor{Apricot!50}}c|}
\hline
\textbf{Costo} & \textbf{Resp.} & \textbf{Amm.} & \textbf{Ana.} & \textbf{Proget.} & \textbf{Progr.} & \textbf{Ver.} & \textbf{Totale} \\
\hline
Costo Orario &  30 \euro &  20 \euro &  25 \euro&  25 \euro &  15 \euro &  15 \euro & /\\
Costo Totale &  120 \euro & 20 \euro & 25 \euro & 375 \euro & 270 \euro & 165 \euro &  975 \euro\\
\hline
\hiderowcolors
\rowcolor{green!50!black!20!white}
\textbf{Saldo fine periodo} & 1140 \euro & 720 \euro & 675 \euro & 2750 \euro & 2085 \euro & 2070 \euro & \cellcolor{yellow!50} 9.440 \euro  \\
\hline
\end{tabular}
\caption{Riepilogo costi e budget rimanente - Quarto Sprint}
\end{table}


\subsubsection{Aggiornamento del Preventivo a Finire - Sprint 4 (Chiusura RTB)}
Al termine del quarto e ultimo \glslink{sprint}{Sprint} della fase di pre-qualifica, viene rendicontato il \textit{\glslink{preventivo}{Preventivo} a Finire} (EAC) finale per la \textit{\glslink{rtb-requirements-and-technology-baseline}{Requirements and Technology Baseline}} (\glslink{rtb-requirements-and-technology-baseline}{RTB}).

La realizzazione e il collaudo del Proof of Concept (PoC) hanno richiesto un effort tecnico superiore alle stime (complessive +7 ore sui ruoli di Progettazione, Programmazione e Verifica). Questo scostamento, sommato a quelli delle iterazioni precedenti, porta la Varianza a Finire complessiva a -300,00 \euro. Trattandosi della chiusura di una macro-fase, il team ha assorbito tale varianza all'interno del margine di \glslink{rischio}{rischio} concordato; la futura \glslink{pianificazione}{pianificazione} della fase di \textit{\glslink{pb-product-baseline}{Product Baseline}} (\glslink{pb-product-baseline}{PB}) sarà calibrata su metriche più conservative per garantire l'assoluto rispetto del tetto di spesa.

\begin{table}[H]
    \centering
    \renewcommand{\arraystretch}{1.5}
    \resizebox{\textwidth}{!}{%
    \begin{tabular}{|l|c|c|c|c|c|}
    \hline
    \textbf{Ruolo} & \textbf{Ore Prev. Iniziali} & \textbf{Scostamento Ore} & \textbf{EAC (Ore)} & \textbf{EAC (Costo)} & \textbf{Varianza (VAC)} \\
    \hline
    \glslink{responsabile}{Responsabile} & 56 & +0 & 56 & 1.680,00 \euro & 0,00 \euro\\
    \hline
    \glslink{amministratore}{Amministratore} & 52 & -1 & 51 & 1.020,00 \euro & +20,00 \euro\\
    \hline
    \glslink{analista}{Analista} & 59 & +6 & 65 & 1.625,00 \euro & -150,00 \euro\\
    \hline
    \glslink{progettista}{Progettista} & 143 & +2 & 145 & 3.625,00 \euro & -50,00 \euro \\
    \hline
    \glslink{programmatore}{Programmatore} & 160 & +4 & 164 & 2.460,00 \euro & -60,00 \euro \\
    \hline
    \glslink{verificatore}{Verificatore} & 160 & +4 & 164 & 2.460,00 \euro & -60,00 \euro\\
    \hline
    \rowcolor{green!50!black!20!white}
    \textbf{Totale} & \textbf{630} & \textbf{+15} & \textbf{645} & \textbf{12.870,00 \euro} & \textbf{-300,00 \euro} \\
    \hline
    \end{tabular}%
    }
    \caption{Preventivo a Finire e Varianza cumulativa finale per la fase RTB}
    \label{tab:eac_sprint4}
\end{table}

\subsection{Rendiconto ore rimanenti}
Come si può osservare nella tabella sottostante, il documento riporta il dettaglio delle ore rimanenti a disposizione. I dati sono suddivisi per ogni singolo componente del gruppo e per le rispettive mansioni, mostrando il monte ore residuo calcolato al termine dei quattro \glslink{sprint}{sprint} di progetto.

\begin{table}[htbp]
\centering
\renewcommand{\arraystretch}{1.5} 
\begin{tabularx}{\textwidth}{|X|c|c|c|c|c|c|c|}
\hline
\textbf{Ore rimanenti} & \textbf{Resp.} & \textbf{Amm.} & \textbf{Ana.} & \textbf{Proget} & \textbf{Progr.} & \textbf{Ver.} & \textbf{Totale} \\
\hline
A. Gastaldello & 3 & 7 & 0 & 14 & 17 & 23 & 64 \\
\hline
E. Bellet & 4 & 1 & 5 & 19 & 23 & 17 & 69 \\
\hline
A.Oloni & 4 & 6 & 7 & 15 & 19 & 23 & 74 \\
\hline
D. Belluz & 8 & 0 & 7 & 16 & 23 & 19 & 73 \\
\hline
A. Maule & 3 & 8 & 0 & 18 & 19 & 21 & 69 \\
\hline
L. Soligo & 8 & 7 & 5 & 10 & 15 & 23 & 68 \\
\hline
V. Y. K. Fernandes & 8 & 8 & 2 & 18 & 23 & 12 & 71 \\
\hline
\rowcolor{gray!20}
\textbf{TOTALE} & \textbf{38} & \textbf{37} & \textbf{26} & \textbf{110} & \textbf{139} & \textbf{138} & \textbf{488} \\
\hline
\end{tabularx}
\caption{Rendiconto - Ore rimanenti}
\end{table}

\end{document}
